\documentclass[11pt,reqno]{amsart}
\usepackage{amssymb,mathtools}
\usepackage{enumitem}
\usepackage{microtype}
\usepackage{geometry}
\geometry{}
\usepackage[hidelinks]{hyperref}
\let\affil\address

\newtheorem{theorem}{Theorem}[section]
\newtheorem{lemma}[theorem]{Lemma}
\newtheorem{proposition}[theorem]{Proposition}
\newtheorem{corollary}[theorem]{Corollary}
\theoremstyle{definition}
\newtheorem{definition}[theorem]{Definition}
\theoremstyle{remark}
\newtheorem{remark}[theorem]{Remark}
\numberwithin{equation}{section}

\newcommand{\R}{\mathbb R}
\newcommand{\E}{\mathbb E}
\newcommand{\Pp}{\mathbb P}
\newcommand{\cF}{\mathfrak F}
\newcommand{\Lip}{\operatorname{Lip}}
\newcommand{\sgn}{\operatorname{sgn}}

\title[Empirical Wasserstein-One Matching]
{Deterministic Selection in Empirical Wasserstein-One Matching}
\author{Wentao Lyu}
\affil{Shanghai Center for Mathematical Sciences,
Fudan University}
\subjclass[2020]{Primary 49Q22; Secondary 60D05, 60G55}
\keywords{Optimal transport, empirical matching, Palm measures}
\thanks{The Lean formalization is available at
\url{https://github.com/LvDAO/OT-Selector}.}
\date{}

\begin{document}

\begin{abstract}
We match independent uniform samples from the squares \([0,1]^2\) and
\([2,3]\times[0,1]\) by minimizing total Euclidean distance. The transport
problem between the two uniform distributions has infinitely many optimal
plans. We prove that, as the sample size grows, the empirical optimal plans
converge weakly in probability to a single deterministic plan. This plan
preserves height and has the same joint distribution of paired horizontal
coordinates at every height. We prove convergence by establishing uniqueness
of the stationary local matching obtained after rescaling and identifying
the selected horizontal coupling from its source-Palm distribution.
\end{abstract}

\maketitle

\section{Introduction}
\label{sec:introduction}

Optimal transport with the Euclidean distance cost can have many
minimizers, even for compactly supported measures with densities. We ask
whether independent sampling selects a definite transport plan. For samples
from the uniform distributions on \([0,1]^2\) and
\([2,3]\times[0,1]\), this question is listed as an open problem in the
survey of del Barrio, Gonz\'alez-Sanz, Loubes, and
Rodr\'iguez-V\'itores~\cite[Section~7, Problem~4]{delBarrioEtAl2025}.

We prove that the optimal empirical matchings converge weakly in probability
to a deterministic plan as the sample size grows. We characterize this plan
through the source-Palm distribution of the unique stationary matching
obtained at the scale of nearby sample points.

\subsection{Model and main result}

For each positive integer \(n\), let
\[
X_i=(U_i,T_i)\in[0,1]^2,\qquad
Y_j=(2+V_j,S_j)\in[2,3]\times[0,1],
\qquad 1\le i,j\le n,
\]
where the \(4n\) coordinates are independent and uniform on \([0,1]\).
Let \(\sigma_n\) minimize
\begin{equation}
\frac1n\sum_{i=1}^n
\sqrt{(2+V_{\sigma_n(i)}-U_i)^2+(S_{\sigma_n(i)}-T_i)^2}
\label{eq:empirical-objective}
\end{equation}
over permutations of \(\{1,\ldots,n\}\). A fixed ordering of the
permutations gives a measurable choice whenever there is a tie.
In the coordinate order \((t,u,s,v)\), define
\begin{equation}
\pi_n:=\frac1n\sum_{i=1}^n
\delta_{(T_i,U_i,S_{\sigma_n(i)},V_{\sigma_n(i)})}.
\label{eq:empirical-plan}
\end{equation}
The physical transport plan is its pushforward by
\((t,u,s,v)\mapsto((u,t),(2+v,s))\).
For a coupling \(\kappa\) of two uniform \([0,1]\) laws, write
\[
\Pi_\kappa(dt,du,ds,dv):=dt\,\delta_t(ds)\,\kappa(du,dv).
\]
Thus \(\Pi_\kappa\) preserves height, and its paired horizontal marks
have law \(\kappa\), independently of height.

\begin{theorem}[Deterministic empirical selection]
\label{thm:main-prove}
There is a deterministic coupling \(\kappa_\infty\) of two uniform
\([0,1]\) laws such that
\begin{equation}
\pi_n\xrightarrow{\Pp}\Pi_{\kappa_\infty}
\quad\text{in the weak topology on probability measures on }[0,1]^4.
\label{eq:main-selector}
\end{equation}
Equivalently, for every continuous \(F:[0,1]^4\to\R\),
\[
\frac1n\sum_{i=1}^n
F(T_i,U_i,S_{\sigma_n(i)},V_{\sigma_n(i)})
\xrightarrow{\Pp}
\int_0^1\int_{[0,1]^2}F(t,u,t,v)\,\kappa_\infty(du,dv)\,dt.
\]
\end{theorem}

The theorem uses the iid law at each sample size; samples of different
sizes may have arbitrary dependence. The selected coupling is the
source-Palm distribution of paired horizontal marks in the stationary
microscopic matching. This is the joint law of the horizontal coordinates
seen from a typical source point and its partner.

We can describe all optimal plans for the limiting problem. For
\(x=(u,t)\) and \(y=(2+v,s)\),
\[
|x-y|\ge 2+v-u,
\]
with equality precisely when \(s=t\). Every admissible continuum
coupling integrates the right-hand side to \(2\), and a
height-preserving coupling attains this value. Consequently the optimal
plans are exactly
\begin{equation}
dt\,\delta_t(ds)\,\kappa_t(du,dv),
\label{eq:optimal-face}
\end{equation}
where \(t\mapsto\kappa_t\) is a measurable family of couplings of two
uniform laws, up to equality for almost every \(t\). This description
follows by disintegrating over the common height. Theorem
\ref{thm:main-prove} therefore identifies both a deterministic
horizontal coupling and its constancy across macroscopic heights.

\subsection{Selection and stationary matching}

Secondary variational selection for the distance cost minimizes an
additional energy over the continuum optimal set; its relationship with
transport rays and variational asymptotic expansions is developed by
Ambrosio and Pratelli~\cite{AmbrosioPratelli}. In our model, quadratic
secondary selection is explicit: on the optimal set,
\[
\int |x-y|^2\,d\pi
=4+\int (v-u)^2\,d\pi.
\]
It therefore selects the horizontal translation \(v=u\).
The approximation scheme matters for selection. Gw\'{o}\'zd\'z~\cite{Gwozdz}
constructs weakly convergent absolutely continuous marginals whose
ray-monotone selectors converge to a different optimal plan from the
selector of the limiting marginals. Our result establishes selection for the
iid empirical approximation in \eqref{eq:empirical-objective}.

The local matching problem connects the proof to invariant minimal
matchings. Holroyd, Janson and W\"astlund~\cite{HolroydJansonWastlund}
study point-process matchings that minimize cost under every finite
rematching, including their uniqueness and translation-invariant
realizations. Here the points carry horizontal marks, and a feedback
change of coordinates gives a quadratic local cost. We prove uniqueness
for stationary locally minimizing matchings with zero relative flux and
finite source-Palm second moment in this environment.

Optimal transport between stationary random measures provides another
point of comparison. Huesmann~\cite{Huesmann} studies mean transport cost
per unit volume and approximation by problems on bounded regions, with
uniqueness for an absolutely continuous source. We also pass from finite
problems to a stationary matching. For the atomic source measure in our
setting, regeneration and stationary mass transport establish uniqueness
within the locally minimizing class.

The selection result also has a stochastic homogenization interpretation.
The continuum cost determines the optimal set, while the rescaled sample
geometry determines which plan is selected. Variational asymptotic
expansions study this type of selection by resolving finer scales within
a minimizing set~\cite{AmbrosioPratelli}.

In the present proof, subtracting terms that depend only on a source or
a target exposes a quadratic local cost. Optimality passes to the local
matching limits, and their unique distribution of paired marks determines
the transport plan at every macroscopic height.

\subsection{Proof strategy}

Fix an auxiliary parameter \(0<\lambda<1/8\). A predictable feedback
graph compensates the source--target counting imbalance on the
microscopic rank scale. Section~\ref{sec:feedback} constructs this
graph and controls its current and a comparison matching.
Section~\ref{sec:localization} transfers the comparison estimate to the
exact Euclidean optimizer by a row--column cost gauge, obtaining a
quadratic bound on feedback edge lengths.

Section~\ref{sec:environment} identifies the stationary marked
environment through its martingale problem. Section~\ref{sec:compactness}
then shows that every selected local matching limit is perfect,
locally minimizing, of zero relative flux, and has finite source-Palm
second moment. Its source-Palm law means the law seen from a typical
source event; the paired source and target marks give the candidate
horizontal coupling.

Sections~\ref{sec:kinetic} and~\ref{sec:regenerative-endgame} prove
uniqueness of this stationary matching. A long zero-current hold,
together with a deterministic energy iteration, produces a common
empty cut in any joining of two candidate matchings. Stationary mass
transport and the absence of finite alternating ties then force the
matchings to coincide. Finally, Section~\ref{sec:reconstruction} uses
vertical concentration and blockwise Palm readout to reconstruct
\(\pi_n\) and prove Theorem~\ref{thm:main-prove}.

The feedback graph is built at a fixed outer error tolerance
\(\varepsilon>0\), and its estimates are derived before intersecting
the required high-probability events. The limit
\(\varepsilon\downarrow0\) is taken only after returning to the
original empirical quantities. This order of limits is used throughout.

\section{Finite feedback estimates}
\label{sec:feedback}

We construct a graph that keeps the source and target crossing counts
close.  Matching the crossings in order will give a comparison matching
with controlled squared edge lengths.

Set
\[
a_i:=1-U_i\in[0,1],\qquad b_j:=1+V_j\in[1,2],\qquad
D_{ij}:=a_i+b_j\in[1,3],
\]
and use the rank-scale phase lines
\begin{equation}
\widehat A_i(p):=nT_i+a_ip,\qquad
\widehat B_j(p):=nS_j-b_jp.
\label{eq:1.1}
\end{equation}
Fix once and for all
\begin{equation}
0<\lambda<\frac18,
\qquad
0<\eta<\frac{\lambda}{4(2+\lambda)}.
\label{eq:1.2}
\end{equation}

\paragraph{Auxiliary graphs and outer limits.}
For every \(\varepsilon>0\), choose a fixed \(A=A(\varepsilon)\), set
\(B_n=A\sqrt n\), and construct the auxiliary graph
\(g_{n,\varepsilon}\).  There is a good event
\(\mathcal G_{n,\varepsilon}\) such that
\begin{equation}
\liminf_{n\to\infty}\Pp(\mathcal G_{n,\varepsilon})
\ge 1-\varepsilon,
\label{eq:1.2a}
\end{equation}
and all graph-dependent estimates below hold on
\(\mathcal G_{n,\varepsilon}\), with constants allowed to depend on
\(\varepsilon\) but not on \(n\).  To reduce notation, the superscript
\(\varepsilon\) is suppressed inside a proof carried out at fixed
\(\varepsilon\).

\paragraph{Allocation of error probabilities.}
Whenever a proof uses finitely many high-probability events, we
first apply the corresponding estimates with tolerances
\(\varepsilon/m\), where \(m\) is a fixed deterministic integer, and
then intersect the resulting events.  We relabel the total tolerance as
\(\varepsilon\).  Throughout, ``on the good event'' refers to this final
intersection.  We first prove the stopped conditional estimates with respect
to the filtration at the stated stopping coordinate, and then form the
intersection.

For
\begin{equation}
\mathsf J_n(p,x)
:=
\#\{i:\widehat A_i(p)\le x\}
-
\#\{j:\widehat B_j(p)\le x\},
\label{eq:1.3}
\end{equation}
the unclipped half-space volumes in the bulk are
\begin{equation}
V^-(p,x)=x-\frac p2,\qquad
V^+(p,x)=x+\frac{3p}{2},
\label{eq:1.4}
\end{equation}
and we write
\[
\Delta_n^-(p,x)
=
\#\{i:\widehat A_i(p)\le x\}-V^-(p,x),
\]
\begin{equation}
\Delta_n^+(p,x)
=
\#\{j:\widehat B_j(p)\le x\}-V^+(p,x).
\label{eq:1.5}
\end{equation}
Thus, whenever \eqref{eq:1.4} is unclipped,
\begin{equation}
\boxed{
\mathsf J_n(p,x)
=
\Delta_n^-(p,x)-\Delta_n^+(p,x)-2p .
}
\label{eq:1.6}
\end{equation}

\subsection{Predictable deferred decisions and adaptive rates}

Let
\[
\mathcal X_n=[0,n]\times[0,1],
\qquad
P_i^-=(nT_i,U_i),\qquad
P_j^+=(nS_j,V_j).
\]

\begin{lemma}[Predictable-frontier deferred decisions]
\label{lem:deferred}
Let \(p=g(x)\) be a left-to-right, piecewise \(C^1\), predictable graph
with \(\Lip(g)\le\lambda<1/2\).  Define the swept sets
\[
H_x^-(g)
=
\{(t,u)\in\mathcal X_n:t+g(x)(1-u)\le x\},
\]
\begin{equation}
H_x^+(g)
=
\{(s,v)\in\mathcal X_n:s-g(x)(1+v)\le x\}.
\label{eq:1.7}
\end{equation}
Let \(\mathcal F_x^{g,0}\) be generated by the graph up to \(x\) and by all
labeled source and target points whose trajectories have crossed it by
coordinate \(x\).  In stopping-coordinate statements we use its usual
augmentation
\[
\mathcal F_x^g
:=
\bigcap_{y>x}
\bigl(\mathcal F_y^{g,0}\vee\mathcal N\bigr)
=
\bigcap_{q>x,\ q\in\mathbb Q}
\bigl(\mathcal F_q^{g,0}\vee\mathcal N\bigr),
\]
where \(\mathcal N\) is the family of ambient terminal null events.
Conditional on \(\mathcal F_x^g\), the unrevealed
source labels are independent uniforms on
\(\mathcal X_n\setminus H_x^-(g)\), the unrevealed target labels are
independent uniforms on \(\mathcal X_n\setminus H_x^+(g)\), and the two
families are conditionally independent.  The same assertion holds at every
\(\mathcal F^g\)-stopping coordinate \(\tau\) with \(0\le\tau\le n\) almost
surely.
\end{lemma}

\begin{proof}
Spacelikeness makes both swept families increasing.  Indeed, for \(y>x\),
\[
[y-g(y)(1-u)]-[x-g(x)(1-u)]
\ge(1-\lambda)(y-x)>0,
\]
and
\begin{equation}
[y+g(y)(1+v)]-[x+g(x)(1+v)]
\ge(1-2\lambda)(y-x)>0.
\label{eq:1.8}
\end{equation}
Suppose at a crossing coordinate \(x_0\) that the unrevealed points have
the stated product-uniform law on the current complements.  Until the
next crossing the continuation of the graph is a deterministic function
of the exposed filtration.  If \(C_x\) is either decreasing complement and
\(Z\) is one unrevealed label, then
\[
\Pp\{Z\in C_x\mid Z\in C_{x_0},\mathcal F_{x_0}^g\}
=\frac{|C_x|}{|C_{x_0}|},
\]
and conditional on \(Z\in C_x\), \(Z\) is uniform on \(C_x\).
The conditioning factorizes over unrevealed labels, so absence of a
crossing up to \(x\) preserves the product-uniform law.

Simultaneous first crossings have probability zero.  To see this, freeze
the two initial revealed-label sets and a finite history of the previously
exposed labels.  On each such history cell
the graph continuation is a measurable function of the coordinates in that
history.  If two unused labels were to cross next at the same coordinate,
then, after fixing every coordinate except the height of the first label,
spacelikeness makes the crossing equation injective in that height.  Thus
the exceptional height section contains at most one point.  The height law
is nonatomic, and there are only finitely many initial patterns, histories,
steps, and label pairs, so the union of all simultaneous-crossing events is
null.  The same one-coordinate argument shows that a label lies on a fixed
frontier with probability zero.

At the next crossing exactly one label is exposed almost surely.
Condition first on entrance through a boundary slab of thickness \(h\),
reveal that label and its position, and let \(h\downarrow0\); the
remaining labels retain the product-uniform law on the new complement.
There are at most \(2n\) crossings.

To construct the conditional law, fix a finite labeled history \(h\) and
the tuple \(z\) of revealed coordinates.  On that history cell define
\(Q_{h,z,x}\) by fixing
the revealed coordinates at \(z\) and assigning to every unrevealed source
or target label the normalized restriction of Lebesgue measure to its
current strict complement; take the product over the unrevealed labels.
The boundary-slab induction gives, first for a future cylinder \(E\) and a
generator \(D\) of the raw history sigma-field,
\begin{equation}
 \int_{D\cap\{h\}} Q_{h,z(\omega),x}(E)\,d\Pp(\omega)
 =\Pp(D\cap E\cap\{h\}).
\label{eq:1.8a}
\end{equation}
The kernel is measurable in the revealed tuple: all its cylinder values are
finite products of measurable section areas divided by the corresponding
positive complement areas.  (On terminal zero-area cells its value may be
chosen arbitrarily.)  Summing over the finitely many histories and applying
the pi-lambda theorem first in \(D\) and then in \(E\) extends \eqref{eq:1.8a} to the
whole raw history sigma-field and all Borel future events.  Thus the assembled
kernel is a regular conditional law at each deterministic coordinate.
The identity also holds almost surely after adding ambient terminal null events.

The two definitions of \(\mathcal F_x^g\) above agree: monotonicity of the raw
filtration and density of \(\mathbb Q\) make the rational right-neighbourhoods
cofinal among all right-neighbourhoods.  For rational \(x_m\downarrow x\),
frontier nullity and dominated convergence give
\(Q_{x_m}(E)\to Q_x(E)\) first for cylinders, while reverse-martingale
convergence identifies the limit with conditioning on \(\mathcal F_x^g\).
A second monotone-class argument extends this identification to every Borel
\(E\).

For a stopping coordinate \(0\le\tau\le n\), let
\(\tau_m=2^{-m}\lceil2^m\tau\rceil\).  Apply the deterministic-coordinate
conditional identity on the finitely many range slices
\(\{\tau_m=k2^{-m}\}\).  Then
\(\tau_m\downarrow\tau\), right continuity gives
\(\mathcal F_\tau^g=\bigcap_m\mathcal F_{\tau_m}^g\), and the reverse
martingale convergence theorem passes the conditional expectations to
\(\mathcal F_\tau^g\).

Continuity of the graph and the frontier-null
statement above give monotone convergence of the remaining-set indicators
and their normalizing areas.  Applying this argument to finite products of
bounded Borel test functions identifies the asserted product-uniform
conditional law at \(\tau\).
\end{proof}

For fixed \(A>10\), put \(B_n=A\sqrt n\) and
\begin{equation}
\mathcal R_{A,n}
=
\{(p,x): B_n\le x\le n-B_n,\ |p|\le B_n/8\}.
\label{eq:1.9}
\end{equation}
For every \(n\ge1\), \((p,x)\in\mathcal R_{A,n}\), and \(u,v\in[0,1]\),
\[
\frac{7B_n}{8}\le x-p(1-u)\le n-\frac{7B_n}{8},
\qquad
\frac{3B_n}{4}\le x+p(1+v)\le n-\frac{3B_n}{4}.
\]
Thus all source and target sections lie strictly inside the rank strip
before the curve exits \(\mathcal R_{A,n}\).  We use the unclipped area
formulas on this interval.
Let
\[
m_-(p,x)=\min\{V^-(p,x),n-V^-(p,x)\},
\quad
m_+(p,x)=\min\{V^+(p,x),n-V^+(p,x)\}.
\]

\begin{lemma}[Half-space regularity and adaptive crossing rates]
\label{lem:rates}
There is an event \(\mathcal E_{A,n}\) with
\(\Pp(\mathcal E_{A,n})\to1\) on which, uniformly on
\(\mathcal R_{A,n}\),
\begin{equation}
|\Delta_n^\pm(p,x)|
\le
C\bigl(\sqrt{m_\pm(p,x)\log n}+\log n\bigr).
\label{eq:1.10}
\end{equation}
Moreover
\begin{equation}
\sup_{\mathcal R_{A,n}}
(|\Delta_n^-|+|\Delta_n^+|)
=O_{\Pp}(\sqrt n),
\label{eq:1.11}
\end{equation}
and at \(x_R=n-B_n\),
\begin{equation}
\sup_{|p|\le B_n/8}
(|\Delta_n^-(p,x_R)|+|\Delta_n^+(p,x_R)|)
=
O_{\Pp}(\sqrt{B_n}).
\label{eq:1.12}
\end{equation}

Let now \(p=g(x)\) be any predictable curve in \(\mathcal R_{A,n}\)
with \(|g'|\le\lambda\).  Write \(\tau_i^-\) and \(\tau_j^+\) for the
source and target crossing coordinates and define the pre-crossing counts
\[
D_-(x-):=\#\{i:\tau_i^-<x\},
\qquad
D_+(x-):=\#\{j:\tau_j^+<x\},
\]
and the remaining counts
\[
M_-(x-):=n-D_-(x-),
\qquad
M_+(x-):=n-D_+(x-).
\]
Put \(K(x-)=D_-(x-)-D_+(x-)\) and let
\(w_-(x)=g'_-(x)\) be the predictable left slope.  Away from crossing
coordinates, \(K(x-)=\mathsf J_n(g(x),x)\) and \(w_-(x)=g'(x)\).
The predictable \emph{marked} crossing intensities are
\begin{equation}
r_-(x)\,[1-w_-(x)(1-u)]\,du,
\qquad
r_+(x)\,[1+w_-(x)(1+v)]\,dv,
\label{eq:1.13}
\end{equation}
where
\begin{equation}
r_-(x)
=
\frac{M_-(x-)}{n-V^-(g(x),x)},
\qquad
r_+(x)
=
\frac{M_+(x-)}{n-V^+(g(x),x)}.
\label{eq:1.14}
\end{equation}
Consequently the total source and target jump rates are
\begin{equation}
\lambda_-(x)
=
r_-(x)\Bigl(1-\frac {w_-(x)}2\Bigr),
\qquad
\lambda_+(x)
=
r_+(x)\Bigl(1+\frac{3w_-(x)}{2}\Bigr),
\label{eq:1.15}
\end{equation}
and, on \(\mathcal E_{A,n}\),
\begin{equation}
\sup_x|r_\pm(x)-1|=o(1)
\label{eq:1.16}
\end{equation}
uniformly along every such curve before it exits
\(\mathcal R_{A,n}\).
\end{lemma}

\begin{proof}
The sets in \eqref{eq:1.7}, with deterministic \(p,x\), form two finite-VC
classes.  The standard relative VC Bernstein inequality (obtained by
dyadic peeling according to the smaller of a set and its complement,
followed by Sauer's lemma and Bernstein's inequality) gives \eqref{eq:1.10};
the ordinary VC maximal inequality gives \eqref{eq:1.11}.

For \eqref{eq:1.12}, keep \(x=x_R\) fixed.  As \(p\) varies, each complement is
a nested one-parameter family contained in a deterministic set of area
\(O(B_n)\).  If \(Z\) is the number of sample points in the largest
member, then \(Z=O_{\Pp}(B_n)\) and
\(Z-|C|=O_{\Pp}(\sqrt{B_n})\).  Conditional on \(Z\), the nested counts
are an empirical distribution function of \(Z\) iid points; its maximal
fluctuation is \(O_{\Pp}(\sqrt Z)\).  This proves \eqref{eq:1.12}.

For the rates, invoke Lemma~\ref{lem:deferred}.  Conditional on the
explored filtration, the unrevealed source labels are uniform on a region
of area \(n-V^-\).  Advancing the graph by \(dx\) sweeps source area
\[
\int_0^1[1-w(1-u)]\,du\,dx
=
\Bigl(1-\frac w2\Bigr)dx,
\]
and marked \(u\)-slice area
\([1-w(1-u)]\,du\,dx\).  Multiplication by the current unexplored
density gives the first formula in \eqref{eq:1.13} and the first rate in
\eqref{eq:1.15}.  The target swept area is
\[
\int_0^1[1+w(1+v)]\,dv\,dx
=
\Bigl(1+\frac{3w}{2}\Bigr)dx,
\]
giving the second formulas.

To identify the predictable intensities, fix deterministic \(s<t\) before exit from
\(\mathcal R_{A,n}\) and \(B\in\mathcal F_s^g\).  On each of the finitely many
revealed-label histories, Fubini's theorem and the conditional product law give
the exact identities
\[
\E\bigl[\mathbf 1_B(D_-(t)-D_-(s))\bigr]
=
\E\biggl[\mathbf 1_B\int_{(s,t]}\lambda_-(x)\,dx\biggr],
\]
\begin{equation}
\E\bigl[\mathbf 1_B(D_+(t)-D_+(s))\bigr]
=
\E\biggl[\mathbf 1_B\int_{(s,t]}\lambda_+(x)\,dx\biggr].
\label{eq:1.15a}
\end{equation}
The same calculation with bounded Borel mark tests gives the marked versions
of these identities.  Frontier-nullity removes the endpoints, while the
strictly positive remaining areas in \(\mathcal R_{A,n}\) justify every
normalization.  A monotone-class extension from the predictable rectangles
\(B\times(s,t]\), followed by finite linear combinations and bounded
approximation, proves that the counting processes compensated by
\(\int\lambda_\pm\) are martingales.  This establishes the predictable
intensity formulas \eqref{eq:1.13}--\eqref{eq:1.15}.

Finally, on \(\mathcal R_{A,n}\) the
unexplored volumes are at least \(cB_n\); \eqref{eq:1.10} divided by these
volumes gives \eqref{eq:1.16}.
\end{proof}

\subsection{Inward feedback, current energy, and spacelikeness}

Starting at \(x=B_n\) with \(g_n(B_n)=0\), construct the graph recursively
by
\begin{equation}
g_n'(x)=\lambda\,\sgn K_n(x),
\qquad
K_n(x):=\mathsf J_n(g_n(x),x),
\label{eq:1.17}
\end{equation}
between crossings, using the new value of \(K_n\) immediately after a
crossing and \(\sgn0=0\).  Define
\[
\tau_{\rm sh}
=
\inf\{x\ge B_n:|g_n(x)|=B_n/8\},
\]
\[
\tau_{\rm reg}
=
\inf\{x\ge B_n:|r_-(x)-1|\vee|r_+(x)-1|\ge\eta\},
\]
\begin{equation}
\tau_n
=
(n-B_n)\wedge\tau_{\rm sh}\wedge\tau_{\rm reg}.
\label{eq:1.18}
\end{equation}

\begin{lemma}[Lower occupation from predictable rates]
\label{lem:occupation}
Let \(N\) be a simple counting process with predictable intensity
\(\Lambda(x)\), and assume on a stochastic interval
\([\sigma,\sigma+r]\) that
\(0<c_0\le\Lambda\le C_0<\infty\).  Then there are \(a_0,b_0>0\)
depending only on \(c_0,C_0\) such that
\begin{equation}
\Pp\{N(\sigma+r)-N(\sigma)\le a_0r\mid\mathcal F_\sigma\}
\le e^{-b_0r}.
\label{eq:1.19}
\end{equation}
The upper rate bound also gives exponential upper tails for the number
of jumps on length-\(r\) intervals.
\end{lemma}

\begin{proof}
For \(\theta>0\), the exponential counting-process martingale
\[
\exp\left\{-\theta N(x)
+(1-e^{-\theta})\int^x\Lambda(y)\,dy\right\}
\]
and optional sampling between \(\sigma\) and \(\sigma+r\) give
\[
\Pp\{N(\sigma+r)-N(\sigma)\le ar\mid\mathcal F_\sigma\}
\le
\exp\{\theta ar-(1-e^{-\theta})c_0r\}.
\]
Take, for example, \(\theta=1\) and
\(a_0=c_0(1-e^{-1})/2\).  The upper tail is analogous.
\end{proof}

For the endpoint estimate, we separate the failure time from the
deterministic right endpoint.  Put
\begin{equation}
\widehat\tau_n:=\tau_{\rm sh}\wedge\tau_{\rm reg},
\qquad x_R:=n-B_n,
\label{eq:1.18a}
\end{equation}
and define the modified killing time
\begin{equation}
\tau_n^\dagger:=
\begin{cases}
\widehat\tau_n,&\widehat\tau_n<x_R,\\
+\infty,&\widehat\tau_n\ge x_R.
\end{cases}
\label{eq:1.18b}
\end{equation}
Then \(\tau_n^\dagger\) is a stopping coordinate and the process is
alive at \(x_R\) exactly when there has been no premature bulk failure.

\begin{theorem}[Feedback current estimates]
\label{thm:A1}
Before \(\tau_n\), the feedback current has a uniform drift toward zero.
There are deterministic \(c,C>0\) such that, for all large \(n\),
whenever \(K_n(x)\ne0\),
\begin{equation}
\sgn(K_n(x))(\lambda_-(x)-\lambda_+(x))\le-c,
\qquad
c\le\lambda_-(x),\lambda_+(x)\le C.
\label{eq:1.20}
\end{equation}
Consequently
\begin{equation}
\E\int_{B_n}^{\tau_n}K_n(x)^2\,dx
\le
Cn+C\E|K_n(B_n)|^3
=
O(n),
\label{eq:1.21}
\end{equation}
\begin{equation}
\sup_{B_n\le x\le\tau_n}|K_n(x)|
=
O_{\Pp}(n^{1/4}+\log n),
\label{eq:1.22}
\end{equation}
and
\begin{equation}
K_n(n-B_n)\mathbf 1_{\{\tau_n=n-B_n\}}
=
O_{\Pp}(1).
\label{eq:1.23}
\end{equation}
\end{theorem}

\begin{proof}
Before \(\tau_n\), \(1-\eta<r_\pm<1+\eta\).  If \(K_n>0\),
then \(g_n'=\lambda\), and
\[
\lambda_- -\lambda_+
\le
(1+\eta)\Bigl(1-\frac\lambda2\Bigr)
-
(1-\eta)\Bigl(1+\frac{3\lambda}{2}\Bigr)
=
-2\lambda+\eta(2+\lambda)\le-c.
\]
If \(K_n<0\), then \(g_n'=-\lambda\), and
\[
\lambda_- -\lambda_+
\ge
(1-\eta)\Bigl(1+\frac\lambda2\Bigr)
-
(1+\eta)\Bigl(1-\frac{3\lambda}{2}\Bigr)
=
2\lambda-\eta(2-\lambda)\ge c.
\]
The individual rate bounds follow from \(\lambda<1/8\).

For \(V(k)=|k|^3\), the jump generator satisfies
\begin{equation}
\mathcal LV(k)\le-c_1k^2+C_1.
\label{eq:1.24}
\end{equation}
Dynkin's formula up to \(\tau_n\) gives \eqref{eq:1.21}.  At \(x=B_n\) the shear
is zero, so \(K_n(B_n)\) is the difference of two independent
\(\mathrm{Bin}(n,B_n/n)\) variables and
\[
\E|K_n(B_n)|^3=O(B_n^{3/2})=O(n^{3/4})
\]
for fixed \(A\).

For \(W(k)=e^{\theta|k|}\), sufficiently small \(\theta>0\) gives
\begin{equation}
\mathcal LW(k)\le-aW(k)+b.
\label{eq:1.25}
\end{equation}
We use a stopped excursion estimate to control the maximum current.  For
\(R\ge1\), put
\[
 W_R(k):=\exp\{\theta(|k|-R)_+\}.
\]
The same one-step calculation, including the flat region \(|k|\le R\),
gives
\begin{equation}
 \mathcal L W_R(k)\le-aW_R(k)+b
\label{eq:1.25c}
\end{equation}
with the same positive constants after increasing \(b\), if necessary.
Let \(G_R=\{|K_n(B_n)|\le R\}\), which is
\(\mathcal F_{B_n}\)-measurable, and put \(T_n=x_R-B_n\le n\).
Multiplying the stopped Dynkin identity by \(\mathbf1_{G_R}\), using
\(W_R(K_n(B_n))=1\) on \(G_R\), and discarding the nonnegative terminal
term gives
\begin{equation}
 \E\left[\mathbf1_{G_R}
   \int_{B_n}^{\tau_n}W_R(K_n(x-))\,dx\right]
 \le \frac{1+bT_n}{a}.
\label{eq:1.25d}
\end{equation}
For \(q\ge1\), let \(J_{R,q}\) count outward jumps, through the
right-closed stopping coordinate \(\tau_n\), which enter
\(\{|K_n|\ge R+q\}\) from its complement.  On \(G_R\), attainment of
height \(R+q\) forces \(J_{R,q}\ge1\).  Immediately before every such
jump,
\[
 W_R(K_n(x-))\ge e^{\theta(q-1)}.
\]
The two outward directions are disjoint, and their predictable rates are
bounded by a deterministic constant \(C\).  The compensator identity and
\eqref{eq:1.25d} therefore yield
\begin{equation}
\Pp\!\left(G_R\cap
  \left\{\sup_{B_n\le x\le\tau_n}|K_n(x)|\ge R+q\right\}\right)
 \le \frac{C}{a}(1+bT_n)e^{-\theta(q-1)}.
\label{eq:1.25e}
\end{equation}
Together with the preceding cubic-moment estimate and Markov's inequality,
\begin{equation}
 \Pp\!\left\{\sup_{B_n\le x\le\tau_n}|K_n(x)|\ge R+q\right\}
 \le
 \frac{\E|K_n(B_n)|^3}{R^3}
 +\frac{C}{a}(1+bT_n)e^{-\theta(q-1)}.
\label{eq:1.25f}
\end{equation}
Choose \(R=D\sqrt{B_n}\) and
\(q=2\theta^{-1}\log(n+1)\).  The first term in \eqref{eq:1.25f} is
\(O(D^{-3})\), uniformly in \(n\), and the second is \(O(n^{-1})\).
First choosing \(D\) and then letting \(n\to\infty\) proves
\(\sup|K_n|=O_{\Pp}(\sqrt{B_n}+\log n)\), which is \eqref{eq:1.22}.

For \eqref{eq:1.23}, kill the process at \(\tau_n^\dagger\), send the killed
state to a cemetery point \(\partial\), and set \(W(\partial)=0\).
Killing adds only a nonpositive jump to the generator, hence
\[
\widetilde{\mathcal L}W\le-aW+b.
\]
Dynkin's formula and Gronwall give, for \(B_n\le x\le x_R\),
\begin{equation}
\E\!\left[
e^{\theta|K_n(x)|};\,\tau_n^\dagger>x
\right]
\le
e^{-a(x-B_n)}
\E e^{\theta|K_n(B_n)|}
+\frac ba.
\label{eq:1.25a}
\end{equation}
At \(x=B_n\),
\[
K_n(B_n)=X_n-Y_n,\qquad
X_n,Y_n\stackrel{\rm iid}{\sim}
\operatorname{Bin}\!\left(n,\frac{B_n}{n}\right).
\]
Using \(e^{\theta|z|}\le e^{\theta z}+e^{-\theta z}\) and the binomial
mgf,
\begin{equation}
\E e^{\theta|K_n(B_n)|}
\le
2\exp\!\left\{
2B_n(\cosh\theta-1)+O(B_n^2/n)
\right\}
\le e^{C_\theta B_n}.
\label{eq:1.25b}
\end{equation}
At \(x=x_R\), the event \(\{\tau_n^\dagger>x_R\}\) is exactly
\(\{\widehat\tau_n\ge x_R\}\), equivalently
\(\{\tau_n=x_R\}\).  The first term in \eqref{eq:1.25a} is
\[
\exp\{-a(n-2B_n)+C_\theta B_n\}=o(1),
\]
because \(B_n=A\sqrt n\).  Therefore
\[
\Pp\{|K_n(x_R)|>M,\ \tau_n=x_R\}
\le
e^{-\theta M}\left(\frac ba+o(1)\right),
\]
uniformly in \(n\), proving \eqref{eq:1.23}.
\end{proof}

\begin{theorem}[Spacelike shear control]
\label{thm:A2}
For every \(\varepsilon>0\) there is a finite
\(A=A(\varepsilon)>10\) such that, with \(B_n=A\sqrt n\),
\begin{equation}
\Pp\{\tau_n=n-B_n\}\ge1-\varepsilon-o(1).
\label{eq:1.26}
\end{equation}
On this event
\begin{equation}
\Lip(g_n)\le\lambda<\frac12,
\qquad
\sup_{B_n\le x\le n-B_n}|g_n(x)|=O_{\Pp}(\sqrt n),
\label{eq:1.27}
\end{equation}
and at the right junction
\begin{equation}
q_n:=g_n(n-B_n)
=
O_{\Pp}(\sqrt{B_n})
=
O_{\Pp}(n^{1/4}).
\label{eq:1.28}
\end{equation}
Extend \(g_n\) by the constant \(0\) to the left of \(B_n\) and by the
constant \(q_n\) to the right of \(n-B_n\).  The extension is globally
\(\lambda\)-Lipschitz.
\end{theorem}

\begin{proof}
On \(\mathcal E_{A,n}\), \eqref{eq:1.16} is \(<\eta\) for all large \(n\), so
\(\tau_{\rm reg}\) cannot occur before
\(\tau_{\rm sh}\wedge(n-B_n)\).  While the graph remains in the bulk
rectangle, \eqref{eq:1.6} gives
\begin{equation}
2g_n(x)
=
\Delta_n^-(g_n(x),x)-\Delta_n^+(g_n(x),x)-K_n(x).
\label{eq:1.29}
\end{equation}
By \eqref{eq:1.11}, the bridge supremum divided by \(\sqrt n\) is tight, while
\eqref{eq:1.22} is \(o_{\Pp}(\sqrt n)\).  At the first shear hit,
\(2|g_n|=A\sqrt n/4\).  Therefore
\[
\limsup_n\Pp\{\tau_{\rm sh}<n-B_n\}
\le
\Pp\{Z>A/8\}+o(1)
\]
for a tight random variable \(Z\).  Choose \(A\) so the right side is
\(<\varepsilon\), proving \eqref{eq:1.26}.

On survival, \eqref{eq:1.23}, \eqref{eq:1.29}, and the nested endpoint bound \eqref{eq:1.12} give
\[
2|q_n|
\le
|\Delta_n^-(q_n,x_R)|
+
|\Delta_n^+(q_n,x_R)|
+
|K_n(x_R)|
=
O_{\Pp}(\sqrt{B_n}),
\]
which is \eqref{eq:1.28}.  The constant extensions preserve the Lipschitz bound.
\end{proof}

\paragraph{Predictable continuation on all samples.}
To form the marked rooted laws in Section~4, we define a predictable graph
on every sample.  The compactness estimates above will then apply on the
final good event.  Follow
\eqref{eq:1.17} until
\[
\bar\tau_n:=\widehat\tau_n\wedge(n-B_n),
\]
and, if \(\bar\tau_n<n-B_n\), keep the graph height constant from
\(\bar\tau_n\) onward until all source and target trajectories have crossed; use
the same constant extensions at the two outer tails.  Denote this everywhere-defined
predictable spacelike graph by \(\bar g_n\).  On
\(\{\tau_n=n-B_n\}\), it agrees with the graph already used above.

This continuation supplies an unconditioned finite crossing path and its
regular conditional likelihood.  We use \(\bar g_n\) in the globally marked
constructions below and suppress the bar.  On the good event, all previous
estimates apply to this graph.

We next control the boundary contribution to the energy of the increasing matching.

\begin{lemma}[Boundary current energy]
\label{lem:boundary}
Write \(B=B_n\), \(R=n-B\), and \(q_n=g_n(R)\).  On the left,
\begin{equation}
\E\int_0^B\mathsf J_n(0,x)^2\,dx\le CB^2.
\label{eq:1.30}
\end{equation}
For fixed \(L\), let the \(\mathcal F_R\)-measurable localization be
\begin{equation}
\mathcal A_{n,L}:=
\{\tau_n=R,\ |q_n|\le L\sqrt B,\ |q_n|\le B/8,
\ |K_n(R)|\le L,\ B/2\le m_\pm(R)\le2B\}.
\label{eq:1.31}
\end{equation}
Here \(m_\pm(R)\) are the two remaining label counts.  On this event,
for every deterministic \(x\in[R,n+2B]\),
\begin{equation}
\E[\mathsf J_n(q_n,x)^2\mid\mathcal F_R]\le C_LB,
\qquad
\E\left[\left.\int_R^{n+2B}\mathsf J_n(q_n,x)^2\,dx
\right|\mathcal F_R\right]\le C_LB^2.
\label{eq:1.32}
\end{equation}
For every prescribed tolerance, first choose \(A\) using
Theorem~\ref{thm:A2} and then \(L\)
large enough that the failure probability of \(\mathcal A_{n,L}\) is
below that tolerance for all sufficiently large \(n\).
\end{lemma}

\begin{proof}
At zero shear the two counts are independent binomials with equal mean \(x\),
so \(\E\mathsf J_n(0,x)^2\le2x\), proving \eqref{eq:1.30}.
Conditional on \(\mathcal F_R\), Lemma~\ref{lem:deferred} makes the unseen
labels independent uniforms on their remaining complements.  If these
complements have areas \(C_\pm(x)\), put \(\rho_\pm=m_\pm(R)/C_\pm(R)\).
The constant shear gives
\begin{equation}
C_\pm(R)=B+O(|q_n|),\quad
|C_\pm(x)-(n-x)_+|\le2|q_n|,\quad
|\rho_+-\rho_-|\le C\frac{|q_n|+|K_n(R)|}{B}.
\label{eq:1.32a}
\end{equation}
The two remaining populations at \(x\) are conditionally independent
binomials.  Their difference consequently has conditional mean at most
\(C(|q_n|+|K_n(R)|)\) in absolute value and variance at most \(4B\).
This proves the first bound in \eqref{eq:1.32}; conditional Fubini over an interval
of length \(3B\) proves the second.  The localization probability follows
from Theorem~\ref{thm:A2}, endpoint-current tightness, and the remaining-count bounds.
\end{proof}

\begin{lemma}[Right interval density]
\label{lem:right-density}
For fixed \(L\), set
\begin{equation}
H=n-(8L\sqrt B+1),\qquad e_n=60\log n+\tfrac14.
\label{eq:1.33}
\end{equation}
On \(\mathcal A_{n,L}\), except for conditional probability at most
\(16n^{-3}\), every \(R\le s\le t\le H\) satisfies
\begin{equation}
N_\pm((s,t))\ge\tfrac18(t-s)-e_n.
\label{eq:1.34}
\end{equation}
Here \(N_\pm\) count the physical graph intersections of each color.
\end{lemma}

\begin{proof}
The strip \([R,H]\) is unclipped by the constant shear.  Conditional
on \(\mathcal F_R\), each color's unseen intersections have expected
count at least \((b-a)/4\) in every interval \((a,b)\subset[R,H]\):
its swept area is \(b-a\), its remaining area is at most \(2B\), and
its remaining population is at least \(B/2\).
The binomial lower-tail bound therefore gives
\begin{equation}
\Pp\{N_\pm((a,b))<(b-a)/8-u\mid\mathcal F_R\}
\le2e^{-u/12}.
\label{eq:1.35}
\end{equation}
Apply this to integer endpoints in \([0,n]\), to both colors, with
\(u=60\log n\).  The union bound is at most
\(4(n+1)^2n^{-5}\le16n^{-3}\).  Round the endpoints of an arbitrary
\((s,t)\) inward to integers.  The lost length is at most two, costing
\(2/8=1/4\); intervals shorter than two satisfy \eqref{eq:1.34} trivially.
Every unseen crossing in this strip is an actual graph intersection,
so the resulting lower bound also holds for the full graph counts.
\end{proof}

Theorems~\ref{thm:A1} and~\ref{thm:A2} imply, on the joint good event,
\begin{equation}
\int_B^{n-B}K_n(x)^2\,dx=O_{\Pp}(n),
\label{eq:1.36}
\end{equation}
and the constant extensions make the graph globally spacelike.

\subsection{Energy of the increasing graph matching}

Let \(\widehat\xi_i,\widehat\eta_j\) be the unique intersections
\begin{equation}
\widehat A_i(g_n(\widehat\xi_i))=\widehat\xi_i,\qquad
\widehat B_j(g_n(\widehat\eta_j))=\widehat\eta_j.
\label{eq:1.37}
\end{equation}
Uniqueness follows from \(\Lip(g_n)<1/2\).  Match the two intersection
lists increasingly, with permutation \(\tau_n^g\), and set
\begin{equation}
d_i^g=|\widehat\eta_{\tau_n^g(i)}-\widehat\xi_i|.
\label{eq:1.38}
\end{equation}
Their rank discrepancy is exactly the graph current:
\begin{equation}
\#\{i:\widehat\xi_i\le x\}-\#\{j:\widehat\eta_j\le x\}=K_n(x).
\label{eq:1.39}
\end{equation}
Write \(l(e)\le r(e)\) for an edge's two endpoints.

\begin{lemma}[Stopped lower occupation]
\label{lem:stopped-occ}
If a simple counting process \(N\) has predictable intensity
\(c_0\le\Lambda(x)\le C_0\) before a stopping coordinate \(\tau\), then
there are \(a_0,b_0>0\), depending only on \(c_0,C_0\), such that
for every stopping coordinate \(\sigma\) and deterministic \(u\ge0\),
\begin{equation}
\Pp\{N(\sigma+u)-N(\sigma)\le a_0u,\ \tau\ge\sigma+u
\mid\mathcal F_\sigma\}\le e^{-b_0u}.
\label{eq:1.39a}
\end{equation}
\end{lemma}

\begin{proof}
Apply the exponential counting-process martingale as in
Lemma~\ref{lem:occupation}, stopping at \(\tau\).  On the event displayed,
the compensator between \(\sigma\) and \(\sigma+u\) is at least \(c_0u\).
Optional sampling gives the bound with survival inside the probability.
\end{proof}

\begin{lemma}[Stopped graph rank-gap energy]
\label{lem:rankgap}
For two simple point processes started with no unmatched points, suppose
both predictable intensities remain in \([c_0,C_0]\) before \(\tau\).
Their increasing matching satisfies
\begin{equation}
\E\sum_e\bigl(\min\{r(e),\tau\}-l(e)\bigr)_+^2
\le C\E\int_0^\tau(1+K(x)^2)\,dx.
\label{eq:1.40}
\end{equation}
The same bound holds when only edges completed before \(\tau\) are retained.
\end{lemma}

\begin{proof}
Charge an edge at its actual left endpoint \(x\).  A censored gap longer
than \(u\) requires \(x+u<\tau\), and either \(|K(x)|>c_1u\) or
fewer than \(c_2u\) opposite-color points in \((x,x+u]\).  Thus
Lemma~\ref{lem:stopped-occ} gives
\begin{equation}
\Pp\{\text{censored gap}>u\mid\mathcal F_x\}
\le\mathbf1_{\{|K(x)|>c_1u\}}+Ce^{-c_2u}.
\label{eq:1.41}
\end{equation}
Integration against \(2u\,du\) and summation over genuine left endpoints
using the stopped compensators prove \eqref{eq:1.40}.  Censoring endpoints creates
no new arrivals at \(\tau\).
\end{proof}

\paragraph{Energy before the right boundary.}
Stop the zero-shear left process at its first rate exit from \([1/2,2]\),
capped at \(B\).  This stop reaches \(B\) except on an \(o(1)\) event:
\(N_\pm(B)\le2B\) has exponentially high probability and \(B=o(n)\).
Let \(P_n\) be its censored prefix energy, \(I_n\) the sum of squared
suffixes after \(B\) of initial-backlog edges completed by \(\tau_n\), and
\(U_n\) the energy of edges starting after \(B\) and completed by \(\tau_n\).
These are stopped random variables on the original sample space.  Put
\begin{equation}
Z_n^{\rm pre}:=2P_n+2I_n+U_n.
\label{eq:1.41a}
\end{equation}
The left bound \eqref{eq:1.30} and Lemma~\ref{lem:rankgap} give
\(\E P_n\le C(B+B^2)\).

For the initial backlog, condition on
\(\mathcal F_B\): its ranks have backlogs at most
\(1,\ldots,|K_n(B)|\).  The same occupation argument gives a squared
suffix bound \(C(k^2+1)\) for backlog \(k\).  Consequently
\(\E I_n\le C(n+\E|K_n(B)|^3)\).

For bulk edges the endpoint charging
argument in \eqref{eq:1.41}, now starting at their actual arrivals after \(B\),
and Theorem~\ref{thm:A1} give \(\E U_n\le C(n+\E|K_n(B)|^3)\).  Since
\(\E|K_n(B)|^3=O(B^{3/2})\) and \(B=A\sqrt n\),
\begin{equation}
\boxed{\E Z_n^{\rm pre}\le C_A n.}
\label{eq:1.41b}
\end{equation}
On left and bulk survival, every edge completed by \(R\) is controlled
by this majorant: an edge crossing \(B\) is split there using
\((a+b)^2\le2a^2+2b^2\).

\paragraph{Two-cut energy accounting.}
For an increasing edge with \(R<r(e)\le H\), let
\(d_R(e)=r(e)-\max\{R,l(e)\}\).  Under \eqref{eq:1.34}, rank order implies
\(\tfrac18(r(e)-x)-e_n\le|K_n(x)|\) on this suffix.  Integrating and
using Young's inequality yields
\[
d_R(e)^2\le32\int_{\max\{R,l(e)\}}^{r(e)}|K_n(x)|\,dx+256e_n^2.
\]
At each cut the increasing edges all cross in the same direction, and
their number is \(|K_n(x)|\).  If \(N_R\) counts edges with \(r(e)>R\),
summing the preceding inequality gives
\begin{equation}
2\sum_{R<r(e)\le H}d_R(e)^2
\le64\int_R^H K_n(x)^2\,dx+512N_R e_n^2.
\label{eq:1.42a}
\end{equation}
We next bound the left prefixes of edges crossing \(R\).
Choose a fixed \(W\) so large that each color has at least \(L\) arrivals
in \([R-W,R]\), outside a probability smaller than any prescribed tolerance
on bulk survival.  Lemma~\ref{lem:stopped-occ}, applied to the deterministic
subinterval \([R-W,R-W/2]\), proves this assertion.  Since
\(|K_n(R)|\le L\), rank order then bounds the entire crossing-prefix
contribution by
\begin{equation}
2\sum_{l(e)\le R<r(e)}(R-l(e))^2\le2LW^2.
\label{eq:1.42b}
\end{equation}
Splitting at \(B\) and \(R\), and keeping full lengths for \(r(e)>H\), gives
\begin{equation}
\sum_e|e|^2\le Z_n^{\rm pre}+2LW^2+
64\int_R^H K_n(x)^2\,dx+512N_R e_n^2+E_n^{\rm term},
\label{eq:1.42c}
\end{equation}
where \(E_n^{\rm term}:=\sum_{r(e)>H}|e|^2\).  In particular, \eqref{eq:1.42c}
includes the prefixes of edges crossing the first right cut.

\begin{lemma}[Terminal localization]
\label{lem:terminal-local}
Fix \(a>0\).  On \(\mathcal A_{n,L}\) and the right-density event, also
require \(|K_n(H)|\le a\sqrt B\).  For all sufficiently large \(n\),
\begin{equation}
N_{\rm term}\le N_R\le4B,\qquad
E_n^{\rm term}\le4(10L+8a+9)^2B^2.
\label{eq:1.42d}
\end{equation}
The added current event fails on \(\mathcal A_{n,L}\) with probability
at most \(C_L/a^2\).
\end{lemma}

\begin{proof}
Each edge counted by \(N_R\) uses a remaining endpoint at \(R\), so
\(N_R\le m_-(R)+m_+(R)\le4B\).  Put
\(V=8(a\sqrt B+e_n)\).  For large \(n\), \(R\le H-V\) and
\(e_n\le\sqrt B\).  The interval bound \eqref{eq:1.34} gives at least
\(a\sqrt B\ge|K_n(H)|\) points of each color in \((H-V,H)\).
Thus every edge crossing \(H\) starts after \(H-V\); all right endpoints
are at most \(n+2|q_n|\).  Every terminal length is consequently at most
\[
n+2|q_n|-H+V\le(10L+8a+9)\sqrt B.
\]
This proves \eqref{eq:1.42d}.  The failure bound for \(|K_n(H)|\) is Chebyshev's
inequality applied to the pointwise conditional second moment in \eqref{eq:1.32}.
\end{proof}

\begin{theorem}[Energy of the comparison matching]
\label{thm:I3}
For every \(\varepsilon>0\) there are finite
\(A=A(\varepsilon)\) and \(C=C(\varepsilon)\) such that
\begin{equation}
\boxed{\Pp\left\{\sum_{i=1}^n
|\widehat\eta_{\tau_n^g(i)}-\widehat\xi_i|^2\le Cn\right\}
\ge1-\varepsilon-o(1).}
\label{eq:1.42}
\end{equation}
\end{theorem}

\begin{proof}
Allocate fixed fractions of \(\varepsilon\) to Theorem~\ref{thm:A2},
the boundary localization,
and the backward window in \eqref{eq:1.42b}, choosing \(A,L,W\) in that order.
Then choose \(a\) large enough for Lemma~\ref{lem:terminal-local}.
The expectation estimates \eqref{eq:1.41b} and \eqref{eq:1.32}, with
\(\mathcal A_{n,L}\in\mathcal F_R\), allow deterministic \(K_1,K_2\) such that
\begin{equation}
\Pp\{Z_n^{\rm pre}>K_1n\}<\varepsilon/8,\qquad
\Pp\!\left(\mathcal A_{n,L}\cap
\left\{\int_R^{n+2B}K_n(x)^2\,dx>K_2B^2\right\}\right)<\varepsilon/8.
\label{eq:1.42e}
\end{equation}
The left-rate and right-grid failures are \(o(1)\).  On the remaining
events, \eqref{eq:1.42c}--\eqref{eq:1.42e}, \(N_R\le4B\), and \(e_n^2\le B\) give
\begin{equation}
\sum_e|e|^2\le K_1n+2LW^2+
\bigl[64K_2+2048+4(10L+8a+9)^2\bigr]B^2\le Cn.
\label{eq:1.42f}
\end{equation}
Here \(B^2=A^2n\); the terminal \(O(B^2)=O(n)\) bound suffices.
Adding the allocated failure probabilities proves \eqref{eq:1.42}.
\end{proof}

\subsection{Vanishing normalized graph height}

\begin{corollary}[Vanishing normalized graph height]
\label{cor:I4}
For the feedback graph define
\[
\gamma_n(z):=\frac1n g_n(nz).
\]
Then
\begin{equation}
\Lip(\gamma_n)\le\lambda
\label{eq:1.43}
\end{equation}
deterministically, and for every \(\delta>0\),
\begin{equation}
\Pp\left(
\mathcal G_{n,\varepsilon}\cap
\{\|\gamma_n\|_\infty>\delta\}
\right)\longrightarrow0.
\label{eq:1.44}
\end{equation}
\end{corollary}

\begin{proof}
Rescaling preserves the Lipschitz constant.  The graph starts at zero and is
stopped no later than its first shear exit.  Continuity at that stop and constant
extension give the deterministic bound
\[
\|g_n\|_\infty\le B_n/8,
\qquad
\|\gamma_n\|_\infty\le\frac{A}{8\sqrt n}\longrightarrow0.
\]
On the survival event this stopped graph agrees with the horizon graph used
above.  This proves \eqref{eq:1.44}.  The same normalized-height bound holds without
conditioning for the everywhere-defined stopped continuation.
\end{proof}

The Euclidean assignment is equivalently the minimizer of the residual
\begin{equation}
c_{ij}:=D_{ij}L\!\left(\frac{S_j-T_i}{D_{ij}}\right),
\qquad
L(s):=\sqrt{1+s^2}-1,
\label{eq:1.45}
\end{equation}
because the subtracted horizontal term is permutation invariant.

\section{Quadratic localization of the optimizer}
\label{sec:localization}

We now transfer the squared-length bound from the comparison matching to
the Euclidean optimizer.  Subtracting row and column potentials leaves
the optimizer unchanged and makes the cost comparable to squared distance
along the graph.

Put
\[
\xi_i:=\frac{\widehat\xi_i}{n},\qquad
\eta_j:=\frac{\widehat\eta_j}{n}.
\]
Then
\begin{equation}
T_i+a_i\gamma_n(\xi_i)=\xi_i,\qquad
S_j-b_j\gamma_n(\eta_j)=\eta_j.
\label{eq:2.1}
\end{equation}

\subsection{Exact feedback gauge}

For
\[
\Phi_n(x):=\int_0^x L'(\gamma_n(s))\,ds
\]
define row and column potentials
\[
\alpha_i:=a_iL(\gamma_n(\xi_i))-\Phi_n(\xi_i),\qquad
\beta_j:=b_jL(\gamma_n(\eta_j))+\Phi_n(\eta_j),
\]
and the gauged edge cost
\begin{equation}
\Gamma^{(n)}_{ij}:=c_{ij}-\alpha_i-\beta_j.
\label{eq:2.2}
\end{equation}
Since \(\sum_i\alpha_i+\sum_j\beta_j\) is independent of the permutation,
the Euclidean optimizer \(\sigma_n\) also minimizes
\(\sum_i\Gamma^{(n)}_{i,\sigma(i)}\).

For a fixed pair write
\[
d:=\eta_j-\xi_i,\quad
u:=\gamma_n(\xi_i),\quad
v:=\gamma_n(\eta_j),\quad
D:=a_i+b_j,
\quad
q:=\frac{S_j-T_i}{D}.
\]
By \eqref{eq:2.1},
\begin{equation}
q=\frac{d+a_iu+b_jv}{D}.
\label{eq:2.3}
\end{equation}

\begin{lemma}[Peano--hinge representation]
\label{lem:hinge}
Assume \(d>0\).  For \(t\in\R\) set
\begin{equation}
\begin{aligned}
H_{ij}(t):={}&
D(q-t)_+-a_i(u-t)_+-b_j(v-t)_+\\
&-\bigl|\{x\in[\xi_i,\eta_j]:\gamma_n(x)>t\}\bigr|.
\end{aligned}
\label{eq:2.4}
\end{equation}
Then
\begin{equation}
\Gamma^{(n)}_{ij}
=\int_{\R}L''(t)H_{ij}(t)\,dt,
\label{eq:2.5}
\end{equation}
and
\begin{equation}
H_{ij}(t)\ge0\qquad\text{for all }t.
\label{eq:2.6}
\end{equation}
For \(d<0\), set
\[
\xi'=-\xi_i,\qquad \eta'=-\eta_j,\qquad
\widetilde\gamma(x)=-\gamma_n(-x).
\]
Keep the source coefficient \(a_i\) attached to \(\xi'\) and the target
coefficient \(b_j\) attached to \(\eta'\).  Then
\[
d'=-d>0,\qquad u'=-u,\qquad v'=-v,\qquad q'=-q.
\]
The reflected gauged functional equals the original one because \(L\) is
even and \(L'\) is odd.
\end{lemma}

\begin{proof}
The linear functional in \(L\) defined by the right side of \eqref{eq:2.2}
annihilates affine functions.  The standard Peano representation modulo
affine functions therefore gives
\[
\mathcal T(L)=\int_{\R}L''(t)\,\mathcal T((\,\cdot-t\,)_+)\,dt
\]
for this functional \(\mathcal T\).  Evaluating
\(\mathcal T((\,\cdot-t\,)_+)\) with \eqref{eq:2.3} gives exactly \(H_{ij}(t)\),
and hence \eqref{eq:2.5}.

It remains to prove \eqref{eq:2.6}.  Since
\[
|v-u|\le\lambda d,
\]
and \(a_i\le1,\ b_j\le2,\ \lambda<1/2\),
\begin{equation}
q-u=\frac{d+b_j(v-u)}D>0,\qquad
q-v=\frac{d-a_i(v-u)}D>0.
\label{eq:2.7}
\end{equation}
Moreover \(q\) lies above the whole graph segment.
Indeed, if for some \(x=\xi_i+s\) both
\(u+\lambda s>q\) and \(v+\lambda(d-s)>q\), then
\[
d>\frac{(q-u)+(q-v)}{\lambda}.
\]
But, with \(r=(v-u)/d\),
\[
\frac{(q-u)+(q-v)}d
=\frac{2+(b_j-a_i)r}{a_i+b_j}
\ge\frac{2-\lambda}{3}>\lambda,
\]
a contradiction.

We prove nonnegativity by considering the possible positions of \(t\).
If \(t<\min(u,v)\), the first three terms of \eqref{eq:2.4} equal \(d\), while
the last term is at most \(d\).

If \(u\le t<v\), they equal \(d-a_i(t-u)\), whereas the
\(\lambda\)-Lipschitz constraint gives
\[
|\{\gamma_n>t\}|\le d-\frac{t-u}{\lambda},
\]
and \(a_i<1/\lambda\).  The case \(v\le t<u\) is identical, using
\(b_j<1/\lambda\).

If \(\max(u,v)\le t<q\), then
\[
D(q-t)_+=d-a_i(t-u)-b_j(t-v).
\]
If the graph rises above \(t\), it spends at least
\((t-u)/\lambda\) before the first up-crossing and at least
\((t-v)/\lambda\) after the last down-crossing.  Hence
\[
|\{\gamma_n>t\}|
\le
\left[d-\frac{t-u}{\lambda}-\frac{t-v}{\lambda}\right]_+,
\]
and again \(a_i,b_j<1/\lambda\).

Finally, if \(t\ge q\), \eqref{eq:2.7}
and the fact that \(q\) lies above the graph make every term in
\eqref{eq:2.4} zero.  Thus \(H_{ij}\ge0\).
\end{proof}

Let \(Q(s):=s^2/2\), and write \(\Gamma^Q_{ij}\) for the same gauge
functional with \(L\) replaced by \(Q\).  Lemma~\ref{lem:hinge} gives
\begin{equation}
\Gamma^Q_{ij}=\int_\R H_{ij}(t)\,dt.
\label{eq:2.8}
\end{equation}

\begin{lemma}[Quadratic spacelike coercivity]
\label{lem:quad}
For every \(\lambda<1/2\) there are
\(0<c_\lambda<C_\lambda<\infty\) such that
\begin{equation}
c_\lambda d^2\le \Gamma^Q_{ij}\le C_\lambda d^2.
\label{eq:2.9}
\end{equation}
One may take
\begin{equation}
c_\lambda=
\frac{(1-\lambda)(1-2\lambda)}{2(3-4\lambda)}.
\label{eq:2.10}
\end{equation}
\end{lemma}

\begin{proof}
For \(d>0\), put \(\Delta\gamma=v-u\) and \(r=\Delta\gamma/d\).
A direct calculation yields
\begin{equation}
\Gamma^Q_{ij}
=
\frac{(d-a_i\Delta\gamma)(d+b_j\Delta\gamma)}{2D}
+
\int_{\xi_i}^{\eta_j}
\bigl(\operatorname{Chord}_{\xi_i,\eta_j}\gamma_n-\gamma_n\bigr)\,ds.
\label{eq:2.11}
\end{equation}
Among \(\lambda\)-Lipschitz functions with endpoint slope \(r\), the
largest possible area above the chord is
\begin{equation}
\frac{\lambda^2-r^2}{4\lambda}d^2.
\label{eq:2.12}
\end{equation}
Indeed, after subtracting the chord, the deviation \(h\) satisfies
\(h(0)=h(d)=0\) and its a.e. derivative lies in
\([-\lambda-r,\lambda-r]\).  The maximal area is attained by the
bang--bang tent with slope \(\lambda-r\) up to its peak and
\(-(\lambda+r)\) afterwards.  Its peak height is
\((\lambda^2-r^2)d/(2\lambda)\), and its triangular area is exactly
\eqref{eq:2.12}.

Next,
\[
\frac{(1-a r)(1+b r)}{2(a+b)}
\]
is decreasing in both \(a\) and \(b\), so its minimum over
\(a\in[0,1], b\in[1,2]\) occurs at \(a=1,b=2\).
Therefore
\[
\frac{\Gamma^Q_{ij}}{d^2}
\ge
\frac{(1-r)(1+2r)}6
-\frac{\lambda^2-r^2}{4\lambda}.
\]
The right side is minimized on \([-\lambda,\lambda]\) at
\(r=-\lambda/(3-4\lambda)\), giving \eqref{eq:2.10}.  The upper bound follows
immediately from \eqref{eq:2.11}, \(|r|\le\lambda\), and
\(|\int(\operatorname{Chord}\gamma_n-\gamma_n)|\le \lambda d^2/2\).

For \(d<0\), the same estimate follows from
\(x\mapsto-x\), \(\widetilde\gamma(x)=-\gamma_n(-x)\), and the evenness
of the quadratic cost.
\end{proof}

\begin{lemma}[Euclidean coercivity]
\label{lem:euclid-coerc}
On the event \(\|\gamma_n\|_\infty\le1\),
\begin{equation}
2^{-3/2}c_\lambda d^2
\le
\Gamma^{(n)}_{ij}
\le
C_\lambda d^2
\qquad\text{for every }i,j.
\label{eq:2.13}
\end{equation}
\end{lemma}

\begin{proof}
For the Euclidean \(L\),
\[
L''(t)=(1+t^2)^{-3/2}.
\]
The support of the nonnegative kernel in Lemma~\ref{lem:hinge} lies
between the graph values and \(q\).  Here
\[
|q|=\left|\frac{S_j-T_i}{D}\right|\le1,
\qquad
|u|,|v|\le1.
\]
Hence \(2^{-3/2}\le L''\le1\) on the support.  Combining
\eqref{eq:2.5}, \eqref{eq:2.8}, and Lemma~\ref{lem:quad} proves \eqref{eq:2.13}.
\end{proof}

\begin{theorem}[Localization of the Euclidean optimizer]
\label{thm:fgl2}
Let \(\sigma_n\) be the exact Euclidean optimal permutation.  Then, for every \(\varepsilon>0\), the corresponding auxiliary
feedback graph satisfies
\begin{equation}
\boxed{
\Pp\left\{
\sum_{i=1}^n
\bigl|\widehat\eta_{\sigma_n(i)}-\widehat\xi_i\bigr|^2
\le C_\varepsilon n
\right\}
\ge1-\varepsilon-o(1).
}
\label{eq:2.14}
\end{equation}
\end{theorem}

\begin{proof}
Fix \(\varepsilon>0\) and allocate fixed fractions of this tolerance to
Theorem~\ref{thm:I3} and Corollary~\ref{cor:I4}.  On the intersection
of their good events,
\(\|\gamma_n\|_\infty\le1\) for all large \(n\), and the optimizer
property together with the row/column gauge gives
\[
\sum_i\Gamma^{(n)}_{i,\sigma_n(i)}
\le
\sum_i\Gamma^{(n)}_{i,\tau_n^g(i)}.
\]
Apply Lemma~\ref{lem:euclid-coerc} and multiply by \(n^2\):
\[
2^{-3/2}c_\lambda
\sum_i|\widehat\eta_{\sigma_n(i)}-\widehat\xi_i|^2
\le
C_\lambda
\sum_i|\widehat\eta_{\tau_n^g(i)}-\widehat\xi_i|^2.
\]
Theorem~\ref{thm:I3} then yields \eqref{eq:2.14} after adding the allocated
failure probabilities.

\end{proof}

\paragraph{Joint high-probability event.}
For each outer tolerance \(\varepsilon\), apply the preceding
estimates with a fixed finite subdivision of \(\varepsilon\), and let
\(\mathcal G_{n,\varepsilon}\) denote their intersection, including the
energy event
\begin{equation}
\frac1n\sum_i
|\widehat\eta_{\sigma_n(i)}-\widehat\xi_i|^2\le C_\varepsilon.
\label{eq:2.17}
\end{equation}
After relabeling the total tolerance,
\begin{equation}
\Pp(\mathcal G_{n,\varepsilon})\ge1-\varepsilon-o(1).
\label{eq:2.18}
\end{equation}
\section{Stationary feedback environment}
\label{sec:environment}

We identify the local environment seen from a uniformly chosen point
in the bulk.  Spatial averaging turns the finite martingale identities
into a stationary marked process.

For fixed \(\varepsilon>0\), choose a deterministic sequence
\(h_{n,\varepsilon}\to\infty\),
\(h_{n,\varepsilon}=o(n)\), with
\(h_{n,\varepsilon}\ge 2B_n\), such that
\begin{equation}
\Pp\!\left(
\mathcal G_{n,\varepsilon}\cap
\left\{\|g_{n,\varepsilon}\|_\infty>
\frac{h_{n,\varepsilon}}4\right\}
\right)\longrightarrow0.
\label{eq:3.0}
\end{equation}
Corollary~\ref{cor:I4} ensures this choice.  Set
\[
I_{n,\varepsilon}:=[h_{n,\varepsilon},n-h_{n,\varepsilon}].
\]
Let \(Z_n\) be uniform on \(I_{n,\varepsilon}\), independently conditional on the finite
sample.  Recenter the feedback scene at \(Z_n\), and normalize the graph
height by
\[
(\theta_{Z_n}g_n)(x):=g_n(Z_n+x)-g_n(Z_n).
\]

At fixed \(\varepsilon\), write \(I_n:=I_{n,\varepsilon}\) and
\(\mathbf Q_n^{\rm env}\) temporarily for
\(\mathbf Q_{n,\varepsilon}^{\rm env}\) inside local calculations.
Define the random empirical environment law
\begin{equation}
\mathbf Q_n^{\rm env}
:=
\frac1{|I_n|}\int_{I_n}\delta_{\theta_x\omega_n^{\rm env}}\,dx.
\label{eq:3.0a}
\end{equation}

\begin{lemma}[Spatial martingale closure]
\label{lem:spatial-mp}
Let \(G\) be a bounded local functional depending only on the rooted
past window \([-R,0]\), let \(f:\mathbb Z\to\mathbb R\) have finite
support, and let \(t>0\).  Write
\[
\Psi_{f,t}(\omega)
=
f(K(t))-f(K(0))-\int_0^t\mathcal Lf(K(s-))\,ds,
\]
where
\[
\mathcal Lf(k)
=b_k[f(k+1)-f(k)]+d_k[f(k-1)-f(k)].
\]
Then, for every \(\delta>0\),
\begin{equation}
\Pp\left(
\mathcal G_{n,\varepsilon}\cap
\left\{\left|\int G\,\Psi_{f,t}\,d\mathbf Q_{n,\varepsilon}^{\rm env}\right|>\delta\right\}
\right)\longrightarrow0.
\label{eq:3.0b}
\end{equation}
The analogous assertions hold for the two marked martingale residuals
with bounded mark tests \(h_-\) and \(h_+\).
\end{lemma}

\begin{proof}
Fix \(\varepsilon>0\), and let \(\tau_{n,\varepsilon}\) be the stopping
coordinate in the feedback construction.  Use the compensated current
martingale stopped at \(\tau_{n,\varepsilon}\), denoted
\(M_{n,\tau}^f\).  Since \(h_{n,\varepsilon}\ge2B_n\), every fixed
past and future observation window rooted in \(I_n\) lies inside the
bulk for all sufficiently large \(n\).  For these roots put
\begin{equation}
X_x:=G(\theta_x\omega_n)
[M_{n,\tau}^f(x+t)-M_{n,\tau}^f(x)].
\label{eq:3.0c}
\end{equation}
If \(x+t\le y\), then \(X_xG(\theta_y\omega_n)\) is measurable
at time \(y\), so the martingale property gives \(\E X_xX_y=0\).
The same holds when \(y+t\le x\).  Bounded \(f,G\) and bounded
stopped rates give \(\E X_x^2\le C_{f,G,t}\) uniformly in \(x,n\).
Ordinary Fubini and Cauchy--Schwarz therefore give
\begin{equation}
\E\left|\frac1{|I_n|}\int X_x\,dx\right|^2
\le\frac{C_{f,G,t}}{|I_n|^2}
\iint_{I_n^2}\mathbf1_{\{|x-y|<t\}}\,dx\,dy
\le\frac{2tC_{f,G,t}}{|I_n|}\longrightarrow0.
\label{eq:3.0d}
\end{equation}
On \(\mathcal G_{n,\varepsilon}\), the stopping coordinate is the right
bulk endpoint, so the stopped increment equals the desired feedback
martingale increment on every admissible rooted window.

The finite generator also differs from \(\mathcal L\) by uniformly
\(o(1)\) on this event.  To obtain this estimate at the fixed stopping tolerance,
put \(q_n=\log n/(13B_n/16)\) and
\(\eta_n=2\cdot33280(\sqrt{q_n}+q_n)\).  Then \(\eta_n\to0\), and
\[
4(33280)^2\log n\le\eta_n^2(13B_n/16),\qquad
2\cdot33280\log n\le\eta_n(13B_n/16).
\]
The same VC good event therefore controls residual densities by
\(|r_\pm-1|\le\eta_n\) throughout the surviving bulk.  A mark test
bounded by \(C\) has drift-integral error at most \(2Ct\eta_n\).
Chebyshev's inequality gives \eqref{eq:3.0b}.

The marked assertions use the same orthogonality argument;
bounded marks and stopped rates bound the second moments of their
fixed-window compensated jump counts.  All martingale estimates are
taken under the original sample law, and then restricted to
\(\mathcal G_{n,\varepsilon}\).
\end{proof}

\begin{lemma}[Marked feedback martingale problem]
\label{lem:marked-mp}
Consider a marked càdlàg process whose current takes values in
\(\mathbb Z\).  For every finitely supported \(f:\mathbb Z\to\mathbb R\)
and bounded continuous \(h_-,h_+:[0,1]\to\mathbb R\), require the processes
\begin{equation}
\begin{aligned}
M_f(t):={}&f(K(t))-f(K(0))\\
&-\int_0^t\Bigl[
b_{K(s-)}(f(K(s-)+1)-f(K(s-)))\\
&\hspace{24mm}
+d_{K(s-)}(f(K(s-)-1)-f(K(s-)))
\Bigr]\,ds
\end{aligned}
\label{eq:3.MP1}
\end{equation}
and
\begin{equation}
M_-(t)
:=
\sum_{\substack{0<s\le t\\ \Delta K(s)=+1}}h_-(U_s)
-
\int_0^t\int_0^1
h_-(u)[1-\lambda\sgn K(s-)(1-u)]\,du\,ds,
\label{eq:3.MP2}
\end{equation}
\begin{equation}
M_+(t)
:=
\sum_{\substack{0<s\le t\\ \Delta K(s)=-1}}h_+(V_s)
-
\int_0^t\int_0^1
h_+(v)[1+\lambda\sgn K(s-)(1+v)]\,dv\,ds
\label{eq:3.MP3}
\end{equation}
are martingales, where
\[
b_k=1-\frac{\lambda}{2}\sgn k,\qquad
d_k=1+\frac{3\lambda}{2}\sgn k.
\]
This marked martingale problem is nonexplosive and uniquely determines
the marked feedback path once the law of \(K(0)\) is specified.
\end{lemma}

\begin{proof}
Let \(\tau\) be the first future event.  Conditional on the entire
initial past and \(K(0)=k\), stop the compensators before \(\tau\).
The survival equation and the two marked identities give, with
\(q_k=b_k+d_k\),
\[
\begin{aligned}
\Pp\{\tau\in dt,\ {-}\text{ event},\ U\in du\mid\mathcal F_0\}
 &=e^{-q_kt}[1-\lambda\sgn(k)(1-u)]\,dt\,du,\\
\Pp\{\tau\in dt,\ {+}\text{ event},\ V\in dv\mid\mathcal F_0\}
 &=e^{-q_kt}[1+\lambda\sgn(k)(1+v)]\,dt\,dv.
\end{aligned}
\]
Optional sampling gives the same conditional kernel after each event.
The bound \(0<2-\lambda\le q_k\le2+\lambda\) excludes explosion and
ensures each next waiting time is finite.  Iterating these conditional
holding, colour and mark kernels determines all finite histories and
hence the entire future path, given the initial current law.
\end{proof}

\begin{lemma}[Deterministic sample selection on good events]
\label{lem:good-sample-extraction}
Fix \(\varepsilon>0\).  Suppose events
\(E_n\subset\mathcal G_{n,\varepsilon}\) have probability at least
\(p>0\) along a subsequence.  One can select samples from these
events and a further increasing subsequence so that all countably many
fixed-window spatial martingale residuals vanish, local count second
moments are uniformly bounded on each fixed window, and every weak
limit has no simultaneous two-colour events.  The root-current
moment bound and every additional finite collection of errors tending
to zero in good-event probability may be imposed simultaneously.
\end{lemma}

\begin{proof}
For integer \(R\ge1\), let \(N_R\) count both colours in \([-R,R]\)
and put \(V_{n,R}=\int N_R^2\,d\mathbf Q_n^{\rm env}\).
Integrating the stopped count estimates over the root gives
\(\E[\mathbf1_{\mathcal G_{n,\varepsilon}}V_{n,R}]
\le C_{\varepsilon,R}\) for each fixed \(R\) and all sufficiently
large \(n\).  Choose \(K_R\) so that
\(C_{\varepsilon,R}/K_R\le p2^{-R-4}\).
Markov's inequality then bounds the probability of each good-event
failure \(V_{n,R}>K_R\) by \(p2^{-R-4}\).

For the collision constraints, cover \([-R,R]\) by overlapping open
windows of length \(O(\delta_{R,l})\).  The two-colour stopped factorial
moment bound gives probability at most \(36\Delta^2\) for two events
in a window of length \(\Delta\).  The expected sum of empirical
collision masses over the grid is consequently at most
\(C_R\delta_{R,l}\).  Require that sum to be at most
\(\sqrt{\delta_{R,l}}\), choosing the meshes to decrease to zero and
satisfy \(\sum_{R,l}C_R\sqrt{\delta_{R,l}}<p/8\).
The same good-event indicator controls the sum over all windows.

For each grid, the condition that every supported continuous test
\(0\le f\le1\) has total two-colour mass at most one is closed in
the vague topology.  Portmanteau passes these grid bounds to a weak
limit.  Letting \(l\to\infty\), and then taking the countable union
over \(R\), excludes both multiplicities and cross-colour collisions.

Enumerate the spatial residuals as \(R_{n,j}\).  At step \(m\), choose
\(n_m\) sufficiently large that the preceding estimates for the
first \(m\) constraints apply and the total failure probability of
\(|R_{n_m,j}|\le1/m\), \(j\le m\), is below \(p/8\).
Impose also the first \(m\) requested vanishing-error bounds.
Their combined loss can be made smaller than \(p\), leaving a sample
in \(E_{n_m}\).  Each step imposes a finite prefix of the constraints,
and every required observation window eventually fits inside the bulk.
For every fixed constraint the resulting sequence satisfies it on a tail.

The root-current bound follows directly from the localization estimate.
For a finite permutation, current equals signed crossing count, so with
\(d_i=|\widehat\eta_{\sigma_n(i)}-\widehat\xi_i|\),
\[
\int |K(0)|\,d\mathbf Q_n^{\rm env}
 \le \frac1{|I_n|}\sum_i d_i
 \le \frac{n+\sum_i d_i^2}{2|I_n|}
 \le \frac{n(1+C_\varepsilon)}{2|I_n|}.
\]
This gives tightness of the integer root current on every selected
sample, directly from \eqref{eq:2.17}.  Together with the local count bounds,
compact marks and the \(\lambda\)-Lipschitz normalized graphs, it
proves tightness of the environment laws.  The local second-moment
bounds give uniform integrability of all fixed-window counts.
\end{proof}

\begin{theorem}[The stationary feedback limit]
\label{thm:env}
Let \(\mathbf Q_{n,\varepsilon}^{\rm env}\) be the empirical environment
law generated by the feedback graph associated with the outer tolerance
\(\varepsilon\).  Then, for every compatible metric \(d_{\rm loc}\) on
local environment laws and every \(\delta>0\),
\begin{equation}
\boxed{
\lim_{\varepsilon\downarrow0}
\limsup_{n\to\infty}
\Pp\left\{
d_{\rm loc}\bigl(
\mathbf Q_{n,\varepsilon}^{\rm env},\cF_\lambda
\bigr)>\delta
\right\}=0.
}
\label{eq:3.0e}
\end{equation}
where \(\cF_\lambda\) is the unique two-sided stationary marked feedback
law described below.  Its current
\(K(x)\in\mathbb Z\) is the birth--death chain
\[
k\to k+1
\quad\text{at rate}\quad
b_k:=1-\frac{\lambda}{2}\sgn k,
\]
\begin{equation}
k\to k-1
\quad\text{at rate}\quad
d_k:=1+\frac{3\lambda}{2}\sgn k.
\label{eq:3.1}
\end{equation}
Conditional on \(K(x-)=k\), a \(+1\) jump carries a source mark with
intensity
\begin{equation}
[1-\lambda\sgn(k)(1-u)]\,du,
\label{eq:3.2}
\end{equation}
and a \(-1\) jump carries a target mark with intensity
\begin{equation}
[1+\lambda\sgn(k)(1+v)]\,dv.
\label{eq:3.3}
\end{equation}
The graph cocycle is
\begin{equation}
g(y)-g(x)=\int_x^y\lambda\sgn K(s)\,ds.
\label{eq:3.4}
\end{equation}
The invariant current law is
\begin{equation}
\varpi_0=\frac{\lambda}{1+\lambda},
\label{eq:3.5}
\end{equation}
\begin{equation}
\varpi_k=
\frac{\varpi_0}{1+3\lambda/2}
\left(\frac{1-\lambda/2}{1+3\lambda/2}\right)^{k-1},
\qquad k\ge1,
\label{eq:3.6}
\end{equation}
\begin{equation}
\varpi_{-k}=
\frac{\varpi_0}{1+\lambda/2}
\left(\frac{1-3\lambda/2}{1+\lambda/2}\right)^{k-1},
\qquad k\ge1.
\label{eq:3.7}
\end{equation}
In particular,
\begin{equation}
\E_\varpi[\sgn K]=0,
\label{eq:3.8}
\end{equation}
and the stationary source and target mark intensities are exactly
\begin{equation}
du,\qquad dv.
\label{eq:3.9}
\end{equation}
\end{theorem}

\begin{proof}
Fix \(\varepsilon>0\).  If good-event convergence to
\(\cF_\lambda\) failed, there would be a subsequence and events
\[
E_n=\mathcal G_{n,\varepsilon}\cap
 \{d_{\rm loc}(\mathbf Q_{n,\varepsilon}^{\rm env},\cF_\lambda)>\delta\}
\]
of probability at least \(p>0\).  Apply
Lemma~\ref{lem:good-sample-extraction} inside these events.  It gives a
deterministic sequence of empirical laws with a weakly
convergent subsequence, denoted \(Q_n\Rightarrow Q\), while retaining
the deviation and all the residual, count-moment and collision bounds.
Here \(\varepsilon,p\) are fixed before any moment cutoff is chosen.

\emph{Fixed cuts and stationarity.}
Write \(L_n=|I_n|\).  Uniform rooting and Fubini give, for either colour,
any fixed \(x\in\R\) and \(r>0\),
\[
\int N([x-r,x+r])\,dQ_n\le 2r\frac n{L_n}.
\]
Take a continuous compactly supported test equal to one on the marked
fibre above \(x\), supported in this interval, and bounded by one.
Its counting integral is a nonnegative continuous vague observable.
Bounded truncation and Portmanteau pass the upper moment bound to
\(Q\).  Since \(n/L_n\) is bounded, letting \(r\downarrow0\) proves
that the fixed cut \(x\) has no event \(Q\)-almost surely.
This conclusion holds separately for each fixed \(x\).

For a fixed shift \(h\), rerooting is continuous when cuts \(0,h\)
are vacant.  The exact finite-root estimate
\[
|Q_n(F\circ\theta_h)-Q_n(F)|
 \le 2\|F\|_\infty |h|/L_n
\]
therefore passes to the limit and proves stationarity of \(Q\).

\emph{The full marked martingale problem.}
Choose a countable algebra of bounded environment-past cylinders,
generated by continuous marked-count tests in rational past windows
and the integer root current.  The observation filtration consists of
the initial current and the subsequently revealed two-colour counts.
Use a countable uniformly dense family of continuous mark tests,
current tests \(\mathbf1_{\{k\}}\), and rational time intervals.
These cylinders depend on the revealed environment past and are
measurable in the original filtration.
The selected samples satisfy all their residual identities asymptotically
by Lemma~\ref{lem:spatial-mp}.

Simplicity and vacancy of deterministic endpoints imply continuity of
each residual observable \(Q\)-almost everywhere.  A fixed residual
is bounded by \(C(1+N_R)\) for some finite \(R\).
The selected bounds \(\int N_R^2\,dQ_n\le K_R\) give uniform
integrability, so its expectation converges and the limiting residual
has expectation zero.

For a fixed rational interval, the residual
defines a finite signed measure on the past sigma-field.  It vanishes
against the generating cylinder algebra, hence against every bounded
past-measurable function by the monotone-class theorem.  Increasing the
past window recovers the entire past.

The error from uniformly approximating a mark test is bounded in \(L^1\) by
the uniform approximation error times \(\E_QN_R+CR\).
Right continuity and the same second-moment bound extend rational times
to all real times.  Thus the marked compensators are martingales;
their \(L^1\) right continuity also permits usual augmentation.

\emph{The invariant current and two-sided points.}
Stationarity and the current indicator identities give the balance
equations for \(p_k=Q\{K(0)=k\}\).
The probability flux \(p_kb_k-p_{k+1}d_{k+1}\) is constant in \(k\).
Since \(p_k\to0\) and the rates are bounded, this constant is zero.
The resulting detailed balance gives \eqref{eq:3.6}--\eqref{eq:3.7}, and normalization
gives \eqref{eq:3.5} and \eqref{eq:3.8}.

Lemma~\ref{lem:marked-mp} now identifies the
joint law of \(K(0)\) and both complete marked configurations on
\((0,\infty)\).  By stationarity, shift any finite window into this
positive half-line and recover its current by signed event counts.
Increasing the windows identifies the joint law of the root current
and both full two-sided configurations with that of \(\cF_\lambda\).

\emph{The original graph.}
The finite recursive curve has right derivative
\(\lambda\sgn K_n\) on every surviving bulk segment, including the
post-event right neighbourhood.  The fundamental theorem of calculus
therefore gives the exact original-graph identity
\[
g_n(B_n+t)-g_n(B_n+s)
 =\int_s^t\lambda\sgn K_n(B_n+u)\,du,
 \qquad 0\le s\le t\le n-2B_n.
\]
Every fixed rooted window is eventually in this range.  For each
rational \(q\), the bounded residual
\[
\min\left\{1,\left|g(q)-\int_0^q\lambda\sgn K(s)\,ds\right|\right\}
\]
is continuous at every scene with vacant root cut.  Indeed, on a
compact interval the current stabilizes locally away from its event
coordinates; a small open neighbourhood of those coordinates has
arbitrarily small Lebesgue measure, while the velocity is bounded by
\(\lambda\).  This proves continuity of the integral observable.

The finite residual is identically zero for all sufficiently large
\(n\); the almost-everywhere continuous mapping theorem therefore
makes it zero \(Q\)-almost surely.  Intersect the resulting full-measure
sets over rational \(q\), then use continuity of both primitives to
obtain \eqref{eq:3.4} for all endpoints.
Thus the normalized graph is a measurable function of the identified
current and point configurations.  The whole law \(Q\) equals
\(\cF_\lambda\), contradicting the retained deviation.

We have proved, for each fixed \(\varepsilon,\delta>0\),
\begin{equation}
\Pp\bigl(\mathcal G_{n,\varepsilon}\cap
 \{d_{\rm loc}(\mathbf Q_{n,\varepsilon}^{\rm env},\cF_\lambda)>\delta\}\bigr)
 \longrightarrow0.
\label{eq:3.9a}
\end{equation}
Adding the probability of \(\mathcal G_{n,\varepsilon}^c\) yields a limiting failure
probability at most \(\varepsilon\).  We then let
\(\varepsilon\downarrow0\).  Averaging the identified marked
intensities and using \eqref{eq:3.8} gives \eqref{eq:3.9}.
\end{proof}

\begin{corollary}[Local limiting feedback cost]
\label{cor:limit-cost}
Normalize \(g(0)=0\).  If a source event is \((\xi,u)\), a target event
is \((\eta,v)\), and
\[
a=1-u,\qquad b=1+v,\qquad D=a+b,\qquad d=\eta-\xi,
\]
then the rescaled finite gauged costs converge uniformly on bounded
rooted configurations to
\begin{equation}
\boxed{
\begin{aligned}
\mathfrak c_\lambda((\xi,u),(\eta,v))
={}&
\frac{d^2}{2D}
-\frac{ab}{2D}\bigl(g(\eta)-g(\xi)\bigr)^2\\
&+
d\,\frac{ag(\xi)+bg(\eta)}D
-\int_\xi^\eta g(s)\,ds .
\end{aligned}}
\label{eq:3.10}
\end{equation}
Moreover
\begin{equation}
c_\lambda' d^2\le \mathfrak c_\lambda\le C_\lambda' d^2.
\label{eq:3.11}
\end{equation}
\end{corollary}

\begin{proof}
On a fixed rank-scale rooted window,
along every good-event subsequential limit, the normalized graph height at the root
tends to zero by Corollary~\ref{cor:I4},
while graph increments and rank displacements are \(O(1)\).

Apply the exact gauge decomposition of Section~\ref{sec:localization} and Taylor-expand only
the curvature coefficient:
\(L''(o(1))\to1\).
After multiplication by \(n^2\), the affine and common-height terms
cancel through the gauge, and the remaining quadratic expression is
exactly \eqref{eq:3.10}.  Uniformity on compact rooted configurations follows from
the uniform continuity of \(L''\) near zero and the \(\lambda\)-Lipschitz
bound.  The coercivity is the rank-scale form of
Lemmas~\ref{lem:quad}--\ref{lem:euclid-coerc}.
\end{proof}

\section{Compactness of selected matchings}
\label{sec:compactness}

Let \(\mathcal G_{\lambda,2}^0\) be the class of all stationary perfect
locally minimizing zero-relative-flux matching laws \(\Sigma\) over
\(\cF_\lambda\) whose source Palm law has finite feedback-edge second
moment,
\begin{equation}
\int d_0^2\,d\Sigma^-<\infty.
\label{eq:4.0E}
\end{equation}
This is exactly the class reached by the compactness theorem below,
because Lemma~\ref{lem:palm-l2} transfers the localization bound \eqref{eq:2.17}.

Throughout the local lemmas in this section, \(\varepsilon>0\) is fixed.
Whenever a good-event failure has probability bounded away from zero,
Lemma~\ref{lem:good-sample-extraction} selects deterministic samples
with the required moment and residual bounds.  The local
limits below are taken along these selected sequences.  We use the same
selection to recover mark averages over the full sample and over
macroscopic blocks.

For \(\Sigma\in\mathcal G_{\lambda,2}^0\), let
\[
\kappa[\Sigma]
:=
\operatorname{Law}_{\text{source Palm}}(U_0,V_{\Sigma(0)}).
\]

\begin{lemma}[Bulk intersection count]
\label{lem:bulk-count}
At fixed \(\varepsilon\), with \(I_n:=I_{n,\varepsilon}\), for every
\(\delta>0\),
\begin{equation}
\Pp\left(
\mathcal G_{n,\varepsilon}\cap
\left\{\frac1n\#\{i:\widehat\xi_i\notin I_n\}>\delta\right\}
\right)\to0,
\label{eq:4.0}
\end{equation}
and the same statement holds for targets.  Hence the source- and
target-intersection intensities converge to one along every
positive-probability good-event subsequence.
\end{lemma}

\begin{proof}
On \(\mathcal G_{n,\varepsilon}\) outside the \(o(1)\) event in \eqref{eq:3.0},
\[
|\widehat\xi_i-nT_i|
=
a_i|g_n(\widehat\xi_i)|
\le \frac{h_{n,\varepsilon}}{4},
\]
and
\[
|\widehat\eta_j-nS_j|
=
b_j|g_n(\widehat\eta_j)|
\le \frac{h_{n,\varepsilon}}{2}.
\]
Hence a source with
\[
T_i\in\left[\frac{2h_{n,\varepsilon}}{n},1-\frac{2h_{n,\varepsilon}}{n}\right]
\]
must satisfy \(\widehat\xi_i\in I_n\), and the analogous statement holds
for targets.  Therefore the excluded source and target counts are bounded
by the numbers of \(T_i\)'s and \(S_j\)'s in physical boundary layers of
total width \(4h_{n,\varepsilon}/n=o(1)\).  Their expectations are \(O(h_{n,\varepsilon})=o(n)\);
Markov's inequality gives \eqref{eq:4.0}.
\end{proof}

Let \(\omega_n\) denote the finite feedback environment together with the
optimal matching \(\sigma_n\).  Define its empirical Lebesgue-rooted scene law
\begin{equation}
\mathbf Q_n
:=
\frac1{|I_n|}
\int_{I_n}\delta_{\theta_x\omega_n}\,dx.
\label{eq:4.1}
\end{equation}
At fixed \(\varepsilon\), this is the abbreviation
\(\mathbf Q_n=\mathbf Q_{n,\varepsilon}\); similarly
\(\omega_n=\omega_{n,\varepsilon}\) and \(I_n=I_{n,\varepsilon}\)
in the local lemmas below.
A rooted scene records the two marked point measures, the marked edge
measure, the graph cocycle normalized by \(g(0)=0\), and the integer
root current \(K(0)\). The three point measures carry their vague
topologies, the graph carries the local uniform topology, and
\(K(0)\in\mathbb Z\) carries the discrete topology. The current at
other coordinates is recovered from \(K(0)\) and the signed counts of
source and target events between the root and that coordinate.

For a source event \(i\), write
\[
d_i:=|\widehat\eta_{\sigma_n(i)}-\widehat\xi_i|.
\]

\begin{lemma}[No loss of partners]
\label{lem:no-loss}
Fix the outer tolerance \(\varepsilon>0\).  For every \(\delta>0\),
\begin{equation}
\lim_{R\to\infty}\limsup_{n\to\infty}
\Pp\left(
\mathcal G_{n,\varepsilon}\cap
\left\{\frac1n\#\{i:d_i>R\}>\delta\right\}
\right)=0.
\label{eq:4.2}
\end{equation}
The same holds when rooted from target endpoints.  Consequently no
partner is lost in any positive-probability good-event subsequential
limit.
\end{lemma}

\begin{proof}
Put
\[
Z_n:=\frac1n\sum_i d_i^2.
\]
By the energy bound \eqref{eq:2.17}, on \(\mathcal G_{n,\varepsilon}\) one has
\(Z_n\le C_\varepsilon\) outside an \(o(1)\) event.  Since
\[
\frac1n\#\{i:d_i>R\}\le\frac{Z_n}{R^2},
\]
\eqref{eq:4.2} follows.  The same edge set, viewed from target endpoints, gives
the target statement.
\end{proof}


\begin{lemma}[Projective Palm readout]
\label{lem:palm}
Let fixed-sample rooted matching laws \(Q_m\), with root intervals
\(A_m=(a_m,b_m]\) and lengths \(L_m\), converge weakly to a stationary
perfect matching law \(Q\) with environment \(\cF_\lambda\).
Suppose the source count in \(A_m\), divided by \(L_m\), tends to one,
and the source count within distance one of its two endpoints,
divided by \(L_m\), tends to zero.  Then for every continuous mark test
\(f:[0,1]^2\to\R\),
\begin{equation}
\frac1{L_m}\sum_{\widehat\xi_i\in A_m}
 f(U_i,V_{\sigma_m(i)})\longrightarrow\int f\,d\kappa[Q].
\label{eq:4.8}
\end{equation}
These two count assumptions follow from Corollary~\ref{cor:I4} and the uniform height
empirical-distribution bound for both \(A_m=I_m\) and fixed
macroscopic blocks \(A_m=n_mJ\Subset(0,n_m)\).
\end{lemma}

\begin{proof}
Choose \(\varphi\ge0\), continuous, supported in \([-1,1]\), with
\(\int\varphi=1\), and set
\(W_m(x)=\int_{A_m}\varphi(x-z)\,dz\).
Fubini gives the exact identity
\[
\int\sum_{(x,u;y,v)\in\mathrm{edges}}
 \varphi(x)f(u,v)\,dQ_m
 =\frac1{L_m}\sum_i W_m(\widehat\xi_i)f(U_i,V_{\sigma_m(i)}).
\]
Since \(0\le W_m\le1\), its difference from the block sum is bounded
by \(\|f\|_\infty\) times the normalized endpoint-strip source count.

First take \(0\le f\le1\).  Multiply the edge test by continuous
compact cutoffs in both endpoint positions, increasing to one.
The corresponding nonnegative continuous vague observables and
Portmanteau, followed by monotone convergence, give
\[
\int f\,d\kappa[Q]\le
\liminf_m\frac1{L_m}\sum_i W_m(\widehat\xi_i)f(U_i,V_{\sigma_m(i)}).
\]
Campbell's formula identifies the limiting integral.  Perfectness
identifies the test \(f=1\) with the source statistic, whose finite
mass tends to one by the two count assumptions.  Applying the same
lower bound to \(1-f\) gives the reverse upper bound.  Shifting and
scaling proves the assertion for every continuous real \(f\).

To verify the two count assumptions, put
\(\zeta_m=\|g_m\|_\infty/n_m\) and
\(\Delta_m=\sup_{t\in[0,1]}|F_m(t)-t|\), where \(F_m\) is the
empirical source-height distribution.  The original intersection
equation gives \(|\widehat\xi_i/n_m-T_i|\le\zeta_m\).
Consequently each block count differs from its physical interval mass
by \(O(n_m\zeta_m+n_m\Delta_m+1)\), and the two endpoint strips
contain at most \(O(n_m\zeta_m+n_m\Delta_m+1)\) sources.
Corollary~\ref{cor:I4} and the VC bound give \(\zeta_m,\Delta_m\to0\) on the selected
samples.  Since \(L_m/n_m\) tends to a positive constant, both
required normalized conclusions follow.
\end{proof}

\begin{lemma}[Inherited source-Palm second moment]
\label{lem:palm-l2}
Let \(\mathbf Q\) be any positive-probability good-event subsequential
scene-law limit at fixed \(\varepsilon\).  Then its source Palm law satisfies
\begin{equation}
\boxed{\int d_0^2\,d\mathbf Q^-<\infty.}
\label{eq:4.12}
\end{equation}
More precisely, the deterministic localization bound on the selected samples gives
\begin{equation}
\int d_0^2\,d\mathbf Q^-\le C_\varepsilon.
\label{eq:4.13}
\end{equation}
\end{lemma}

\begin{proof}
Work along the selected deterministic subsequence, where
\(Z_n=n^{-1}\sum_i d_i^2\le C_\varepsilon\).  Choose a nonnegative
continuous \(\varphi\), supported in \([-1,1]\), with integral one,
and continuous cutoffs \(\chi_R(d)\uparrow1\), supported where
\(d\le R+1\).  The edge statistic
\[
H_R(\omega)=\sum_{(x,u;y,v)\in\mathrm{edges}}
 \varphi(x)|y-x|^2\chi_R(|y-x|)
\]
is a nonnegative continuous vague observable.  Exact Fubini and
\(0\le\int_{I_n}\varphi(\widehat\xi_i-z)\,dz\le1\) give
\(\int H_R\,dQ_n\le(n/|I_n|)Z_n\).
Portmanteau and Campbell's formula imply
\[
\int d_0^2\chi_R(d_0)\,dQ^-
 =\int H_R\,dQ\le C_\varepsilon.
\]
Monotone convergence as \(R\to\infty\) proves \eqref{eq:4.13}.
For roots in \(nJ\), the bound is \(C_\varepsilon/|J|\), still finite.
\end{proof}

\begin{lemma}[Finite crossing number by Campbell mass transport]
\label{lem:finite-crossings}
Let
\[
N_{\rm cross}(0)
=
\#\{\text{selected edges crossing }0\}.
\]
Then
\begin{equation}
\int N_{\rm cross}(0)\,d\mathbf Q
=
\int |d_0|\,d\mathbf Q^-<\infty,
\label{eq:4.16}
\end{equation}
hence \(N_{\rm cross}(0)<\infty\) almost surely.  Moreover
\begin{equation}
\int N_{>R,\rm cross}(0)\,d\mathbf Q
=
\int |d_0|\mathbf1_{\{|d_0|>R\}}\,d\mathbf Q^-
\le
\frac1R\int d_0^2\,d\mathbf Q^-.
\label{eq:4.17}
\end{equation}
\end{lemma}

\begin{proof}
For each selected edge, send unit density along the interval joining its
two feedback endpoints.  Stationary Campbell/mass transport gives
\[
\int N_{\rm cross}(0)\,d\mathbf Q
=
\int |d_0|\,d\mathbf Q^-.
\]
Cauchy--Schwarz and Lemma~\ref{lem:palm-l2} make the right side finite.
Restricting the same transport to \(|d_0|>R\) gives the first equality
in \eqref{eq:4.17}, and
\[
|d_0|\mathbf1_{\{|d_0|>R\}}\le d_0^2/R
\]
gives the estimate.
\end{proof}

\begin{lemma}[Zero-relative-flux closure]
\label{lem:zero-flux}
Every selected good-event local limit of \(\mathbf Q_n\) has zero
relative flux with respect to its feedback current.
\end{lemma}

\begin{proof}
At every non-event cut \(x\), define
\[
C_n(x)
=
\#\{i:\widehat\xi_i\le x<\widehat\eta_{\sigma_n(i)}\}
-
\#\{i:\widehat\eta_{\sigma_n(i)}\le x<\widehat\xi_i\}.
\]
For every finite permutation,
\begin{equation}
\boxed{C_n(x)=K_n(x)}
\label{eq:4.18}
\end{equation}
identically.

For \(R<\infty\), let \(C_n^{(R)}(x)\) count only selected crossing edges
of feedback length at most \(R\).  This is a local observable.  Since
\(C_n=K_n\),
\begin{equation}
\begin{aligned}
\mathbf Q_n\{C_n^{(R)}(0)\neq K_n(0)\}
&\le
\frac1{|I_n|}\sum_{d_i>R}d_i\\
&\le
\frac n{|I_n|}\frac{C_\varepsilon}{R}.
\end{aligned}
\label{eq:4.19}
\end{equation}
Along the selected deterministic subsequence, take \(Q_n\Rightarrow Q\).
For every continuity radius \(R\), vacancy of the root cut and local
edge convergence make the truncated integer-current discrepancy
an almost-everywhere continuous observable.  Passing \eqref{eq:4.19} to the
limit gives
\begin{equation}
\mathbf Q\{C^{(R)}(0)\neq K(0)\}\le\frac{C_\varepsilon}{R}.
\label{eq:4.20}
\end{equation}

By Lemma~\ref{lem:finite-crossings},
\[
N_{>R,\rm cross}(0)\longrightarrow0
\quad\text{in }L^1(\mathbf Q),
\]
and therefore
\[
C^{(R)}(0)\longrightarrow C(0)
\]
in probability (indeed eventually almost surely for the finite realized
set of crossing edges).  Hence
\[
\begin{aligned}
\mathbf Q\{C(0)\neq K(0)\}
&\le
\mathbf Q\{C(0)\neq C^{(R)}(0)\}\\
&\quad+
\mathbf Q\{C^{(R)}(0)\neq K(0)\}
\longrightarrow0
\end{aligned}
\]
along continuity radii \(R\to\infty\).  It follows that
\begin{equation}
C(0)=K(0)\qquad\mathbf Q\text{-a.s.}
\label{eq:4.21}
\end{equation}
At a source event both \(C\) and \(K\) jump by \(+1\), at a target event
both jump by \(-1\), and between events both are constant.  Thus
\(C-K\) is a spatially constant integer, and \eqref{eq:4.21} forces it to be
zero.
\end{proof}

\begin{theorem}[Compactness of locally minimizing matchings]
\label{thm:comp}
Let \(\mathbf Q_{n,\varepsilon}\) be the empirical matching-scene law
built from the feedback graph associated with outer tolerance
\(\varepsilon\).  For every compatible metric \(d_{\rm scene}\) on local
matching-scene laws and every \(\delta>0\),
\begin{equation}
\lim_{\varepsilon\downarrow0}
\limsup_{n\to\infty}
\Pp\left\{
d_{\rm scene}\bigl(
\mathbf Q_{n,\varepsilon},\mathcal G_{\lambda,2}^0
\bigr)>\delta
\right\}=0.
\label{eq:4.C}
\end{equation}
Every good-event subsequential limit \(\mathbf Q\) has the following
properties:

\begin{enumerate}[label=\textup{(\roman*)},leftmargin=2.3em]
\item \(\mathbf Q\) is translation invariant and its environment
marginal is \(\cF_\lambda\);

\item the matching is a perfect matching of the source and target events;

\item the matching is locally minimizing for
\(\mathfrak c_\lambda\) against every finite alternating modification;

\item the matching has zero relative flux;

\item if \(\mathbf Q^{-}\) denotes its source Palm law, then the
pushforward
\[
(\xi=0,u;\ \eta,v=\text{mark of the matched target})
\longmapsto (u,v)
\]
is exactly a subsequential weak limit of
\begin{equation}
\kappa_n:=\frac1n\sum_{i=1}^n
\delta_{(U_i,V_{\sigma_n(i)})}.
\label{eq:4.22}
\end{equation}
In particular both marginals of that projective law are uniform.
\end{enumerate}
\end{theorem}

\begin{proof}
\emph{Tightness and stationarity.}
Fix \(\varepsilon>0\) and work on
\(\mathcal G_{n,\varepsilon}\).  The environment marginal is tight by
Theorem~\ref{thm:env}.  Theorem~\ref{thm:fgl2} gives a bound
\(Z_n\le C_\varepsilon\) outside an additional \(o(1)\) event.  On this
event, Lemma~\ref{lem:no-loss} implies that, in any fixed rooted window,
the probability of seeing a source whose partner lies farther than
\(R\) is \(O_{\varepsilon}(R^{-2})\); the same holds for targets.
More explicitly, for every fixed observation radius \(A\),
\[
\mathbf Q_n\{\exists\hbox{ a source in }[-A,A]\hbox{ with }d_i>R\}
\le
\frac{(2A+o(1))}{|I_n|}\#\{i:d_i>R\}
\le
(2A+o(1))\frac{Z_n}{R^2}.
\]
The same estimate holds for targets.  Together with local finiteness of
the point processes this gives tightness of the matching relation.

Approximate translation invariance
follows directly from the averaging definition \eqref{eq:4.1}:
for a bounded local \(F\) and fixed \(h\),
\[
\left|\int F\circ\theta_h\,d\mathbf Q_n-\int F\,d\mathbf Q_n\right|
\le
2\|F\|_\infty\frac{|h|}{|I_n|}+o(1).
\]
Hence every good-event limit is translation invariant.  Its
environment marginal is \(\cF_\lambda\) by Theorem~\ref{thm:env} on the good event.

\emph{Perfectness.}
Lemma~\ref{lem:no-loss} rules out escape of a partner to infinity under
either source or target Palm rooting.  Since every finite state is a
bijection, the local limit is therefore a bijection as well.

\emph{Palm law, inherited second moment, and zero relative flux.}
Lemma~\ref{lem:palm-l2} transfers the finite-sample second-moment bound directly
through compact edge tests,
and Lemmas~\ref{lem:finite-crossings}--\ref{lem:zero-flux} give a finite
crossing current with zero relative-flux constant.

\emph{Local optimality.}
For each \(m\), let \(F_m\) be the closure of the union of all
rooted optimizer scenes from the selected samples with index at least
\(m\).  Each finite law has mass one on \(F_m\) eventually; closed-set
Portmanteau gives \(Q(F_m)=1\), hence \(Q(\bigcap_mF_m)=1\).

Suppose a simple perfect scene in this intersection admitted a strict
finite improvement by \(\delta>0\).  All involved endpoints lie in
a bounded window.  Simplicity and strictness give a neighbourhood of
that scene in which the same finite modification remains improving.
The compact-window cost convergence of Corollary~\ref{cor:limit-cost}
then excludes every sufficiently late optimizer scene from
that neighbourhood: its gauged cost improvement would also improve
the original Euclidean assignment.  This contradicts membership in
the corresponding \(F_m\).  Thus local minimality holds \(Q\)-almost
surely.

\emph{Projective Palm identification.}
The mass-conservation readout in Lemma~\ref{lem:palm} proves (v).
The first marginal is uniform by the identified source intensity.
For the second, send each target-mark indicator along the matched
edge.  Stationary mass transport identifies its outgoing source-window
mean with the target-window intensity, which is also uniform.

We now prove the outer statement.  If, for some fixed
\(\varepsilon,\delta>0\), the event
\[
E_n:=\mathcal G_{n,\varepsilon}\cap
\left\{d_{\rm scene}(\mathbf Q_{n,\varepsilon},\mathcal G_{\lambda,2}^0)>\delta\right\}
\]
had probability bounded below along a subsequence, apply
Lemma~\ref{lem:good-sample-extraction} inside \(E_n\), retaining the
localization bound and all the fixed-window constraints used above.

The resulting deterministic scene laws are tight.  A convergent subsequence
has the identified environment, perfectness, finite Palm second moment,
local minimality and zero relative flux, so its limit belongs to
\(\mathcal G_{\lambda,2}^0\), contradicting the retained distance
\(>\delta\).
Thus \(\Pp(E_n)\to0\).  Since
\(\Pp(\mathcal G_{n,\varepsilon}^c)\le\varepsilon+o(1)\),
\[
\limsup_{n\to\infty}
\Pp\left\{
d_{\rm scene}\bigl(
\mathbf Q_{n,\varepsilon},\mathcal G_{\lambda,2}^0
\bigr)>\delta
\right\}
\le\varepsilon,
\]
and \(\varepsilon\downarrow0\) gives \eqref{eq:4.C}.
\end{proof}

\section{Kinetic gauge and finite genericity}
\label{sec:kinetic}

Let \((K,g,\mathcal N^-,\mathcal N^+)\) have the stationary feedback law
\(\cF_\lambda\).  The graph \(g\) is specified up to an additive constant by
\begin{equation}
g(y)-g(x)=\int_x^y\lambda\sgn K(s)\,ds.
\label{eq:7.1}
\end{equation}
Write source marks as \(a=1-u\in[0,1]\) and target marks as
\(b=1+v\in[1,2]\).  With respect to the natural left-to-right
filtration, their predictable marked intensities are
\begin{equation}
[1-a\lambda\sgn K(x-)]\,da\,dx,
\qquad
[1+b\lambda\sgn K(x-)]\,db\,dx.
\label{eq:7.2}
\end{equation}

\subsection{Exact kinetic action}

For any representative of the graph cocycle, let
\begin{equation}
X=\xi-a g(\xi),\qquad Y=\eta+b g(\eta),
\qquad
\mathsf h((X,a),(Y,b)):=\frac{(Y-X)^2}{2(a+b)}.
\label{eq:7.3}
\end{equation}
Choose an absolutely continuous primitive \(G' = g\).

\begin{proposition}[Exact kinetic gauge]
\label{prop:kinetic-gauge}
The limiting feedback cost \eqref{eq:3.10} satisfies
\begin{equation}
\boxed{
\mathfrak c_\lambda((\xi,a),(\eta,b))
=\mathsf h((X,a),(Y,b))+R(\xi,a)+C(\eta,b),}
\label{eq:7.4}
\end{equation}
where
\begin{equation}
R(\xi,a)=G(\xi)-\frac a2g(\xi)^2,
\qquad
C(\eta,b)=-G(\eta)-\frac b2g(\eta)^2.
\label{eq:7.5}
\end{equation}
Thus finite alternating comparisons for \(\mathfrak c_\lambda\) are
exactly the finite alternating comparisons for \(\mathsf h\).
\end{proposition}

\begin{proof}
Put \(p=g(\xi)\), \(q=g(\eta)\), \(d=\eta-\xi\), and
\(D=a+b\).  Since \(Y-X=d+ap+bq\), direct expansion gives
\[
\mathfrak c_\lambda-\mathsf h
=-\frac{ab(q-p)^2+(ap+bq)^2}{2D}
-\int_\xi^\eta g(s)\,ds.
\]
The identity
\[
ab(q-p)^2+(ap+bq)^2=D(ap^2+bq^2)
\]
reduces this to \eqref{eq:7.5}.  Row and column terms cancel under every finite
alternating rematching.
\end{proof}

\begin{lemma}[No finite alternating ties]
\label{lem:no-finite-alternating-ties}
Almost surely under the feedback environment law, no finite alternating
cycle has equal costs for its two pairings.
\end{lemma}

\begin{proof}
Condition on the jump coordinates, jump types, and all marks except
one source mark \(a\) on a given cycle.  The feedback path \(g\)
is fixed by the coordinates and jump types.  The two targets competing
at that source have distinct marks \(b_1,b_2\) almost surely.

By the exact kinetic gauge, the cycle cost difference, as a function
of \(a\), is a fixed constant plus the difference of the two rational
edge costs.  The latter have distinct poles at \(-b_1,-b_2\).
The numerator at \(-b_j\) is the square of
\[
\eta_j-\xi+b_j\bigl(g(\eta_j)-g(\xi)\bigr),
\]
which is nonzero: the map
\(x\mapsto x+b_j(g(x)-g(\xi))\) is strictly increasing and
equals \(\xi\) only at \(x=\xi\), while event coordinates
are distinct.  Hence the cost difference is not identically zero.
Its zero set in the admissible interval for \(a\) is finite.

Conditional on the unmarked jump path, the source mark at pre-jump
state \(k\) has density
\[
\frac{1-a\lambda\sgn k}{1-(\lambda/2)\sgn k},\qquad 0<a<1,
\]
and marks at different jumps are conditionally independent.  This
density is bounded and strictly positive, so the tie has probability zero.
Taking the union over the countably many finite label cycles proves the claim.
\end{proof}

\section{Regenerative finite-energy uniqueness}
\label{sec:regenerative-endgame}

The compactness theorem supplies stationary, locally minimizing,
zero-relative-flux matching laws with finite source-Palm second moment.
We prove uniqueness by finding a cut that two such matchings leave empty.
To obtain this cut, we condition on a long remaining zero-current hold at
a fixed spatial origin.  Regeneration describes the conditional
environment, and an energy iteration controls the edges that could cross
the hold.  The finite source-Palm second moment closes the iteration.

For a stationary phase \(\Sigma\), write \(d_e\) for the feedback-coordinate
length of a selected edge \(e\).  For a cut \(x\) and bounded interval \(I\), put
\begin{equation}
T_\Sigma(x):=\#\{e:x\in I_e\},\qquad
S_\Sigma(x):=\sum_{e:x\in I_e}|d_e|,
\label{eq:8.1}
\end{equation}
\begin{equation}
\mathcal H_\Sigma(I)
:=\int_I\bigl(T_\Sigma(x)+S_\Sigma(x)\bigr)\,dx,
\label{eq:8.2}
\end{equation}
where \(I_e\) is the closed interval joining the two feedback endpoints of \(e\).

\begin{lemma}[Mean local crossing energy]
\label{lem:mean-local-energy-prove}
For every \(\Sigma\in\mathcal G_{\lambda,2}^0\), there is a finite constant \(C_\Sigma\) such that
\begin{equation}
\E\mathcal H_\Sigma(I)
=|I|\left(\E_{\Sigma^-}|d_0|+\E_{\Sigma^-}d_0^2\right)
\le C_\Sigma|I|.
\label{eq:8.3}
\end{equation}
\end{lemma}

\begin{proof}
For each selected edge send unit density along its feedback interval.  Stationary
Campbell mass transport gives
\[
\E T_\Sigma(0)=\E_{\Sigma^-}|d_0|.
\]
If an edge of length \(|d|\) sends mass density \(|d|\) along the same interval, the
mass through the origin is \(S_\Sigma(0)\), and Campbell gives
\[
\E S_\Sigma(0)=\E_{\Sigma^-}d_0^2.
\]
Integrate over \(I\) and use Cauchy--Schwarz together with the definition of the finite source-Palm second moment.
\end{proof}

\subsection{Return-cell regeneration}

Let \(\mathcal R\) be the point process of jump coordinates at which the stationary
feedback current jumps from a nonzero state into zero.  Under the Palm law
\(\Pp^{\mathcal R}\), enumerate
\begin{equation}
\cdots<R_{-1}<R_0=0<R_1<R_2<\cdots.
\label{eq:8.4}
\end{equation}

For event-count purposes, the \(j\)-th return cell is the path segment
\((R_j,R_{j+1}]\): it starts immediately after the return jump at \(R_j\), includes the
zero holding time following \(R_j\), all subsequent excursion events, and the return jump
at \(R_{j+1}\).  Let \(L_j=R_{j+1}-R_j\), let \(N_j\) count the
source and target jumps in \((R_j,R_{j+1}]\), and write
\[
Z_j:=\text{its initial zero hold},\qquad E_j:=L_j-Z_j.
\]
The marked excursion consists of all subsequent holding increments, jump
types and marks, starting with the departure from zero and ending at the
next return.  Its duration is \(E_j\).

\begin{lemma}[IID return cells with exponential moments]
\label{lem:return-cells-prove}
Under \(\Pp^{\mathcal R}\), the marked return cells are iid.
The field \((Z_j)_{j\in\mathbb Z}\) consists of iid
\(\operatorname{Exp}(2)\) variables independent of the entire marked
excursion field.  There exists
\(\theta_0>0\) such that
\begin{equation}
\E^{\mathcal R}e^{\theta_0N_0}<\infty,
\qquad
\E^{\mathcal R}e^{\theta_0L_0}<\infty.
\label{eq:8.5}
\end{equation}
Every return cell contains equally many source and target events.
\end{lemma}

\begin{proof}
At state zero the total jump rate is two, so its waiting time is
exponential of rate two and independent of the departure type and mark.
The strong Markov property at the departure makes the subsequent marked
excursion independent of that waiting time.

After leaving zero, the sign of the current cannot change before returning
to zero.  On the positive half-line the embedded jump chain moves away from
zero with probability
\[
p_+
=\frac{1-\lambda/2}{(1-\lambda/2)+(1+3\lambda/2)}
=\frac{1-\lambda/2}{2+\lambda}<\frac12.
\]
On the negative half-line it moves away from zero with probability
\[
p_-
=\frac{1-3\lambda/2}{(1+\lambda/2)+(1-3\lambda/2)}
=\frac{1-3\lambda/2}{2-\lambda}<\frac12.
\]
Thus the number of jumps in an excursion from zero back to zero has an exponential
moment.  The holding rate is always one of \(2,2+\lambda,2-\lambda\), hence is bounded
below by \(2-\lambda>0\).  Conditional on the jump count, the cell length is dominated
by a sum of that many iid exponentials with rate \(2-\lambda\), plus the initial zero
holding time.  Choosing \(\theta_0>0\) sufficiently small gives the second exponential
moment in \eqref{eq:8.5}.

The strong Markov property at each jump into zero gives independence and identical
distribution of the forward cells.  The stationary Palm construction gives the same
law in both directions.

Immediately after \(R_j\) the current is zero, and after the included return jump at
\(R_{j+1}\) it is zero again.  Hence the net number of \(+1\) minus \(-1\) jumps in
\((R_j,R_{j+1}]\) is zero, exactly equality of the source and target counts.
\end{proof}

For a return cell, match its source events increasingly to its target events
in feedback coordinate, and let \(W_j\) be the resulting feedback cost.

\begin{lemma}[Comparison rewards excluding the zero hold]
\label{lem:idle-free-reward}
There are iid nonnegative random variables
\[
V_j=C_\lambda N_jE_j^2
\]
which dominate the cellwise increasing comparison rewards \(W_j\),
are independent of the entire field \((Z_j)_{j\in\mathbb Z}\), and satisfy
\begin{equation}
\Pp^{\mathcal R}\{V_j>t\}\le Ce^{-ct^{1/3}}.
\label{eq:8.6}
\end{equation}
\end{lemma}

\begin{proof}
The cell convention is \((R_j,R_{j+1}]\), so the return event at
\(R_j\) is not one of the events matched inside cell \(j\).  Every
matched vertex in that cell lies in
\([R_j+Z_j,R_{j+1}]\).  Each comparison edge therefore has length
at most \(E_j\).  The upper bound in \eqref{eq:3.11} gives
\(W_j\le C_\lambda N_jE_j^2\).  The variables on the right are
functions only of the marked excursions.  They are iid and independent
of all zero holds.

Both \(N_j\) and \(E_j\) have exponential
moments by the proof of Lemma~\ref{lem:return-cells-prove}.  If
\(N_jE_j^2>t/C_\lambda\), either \(N_j\) or \(E_j\) exceeds
a constant multiple of \(t^{1/3}\), proving \eqref{eq:8.6}.
\end{proof}

\begin{lemma}[The zero hold seen from a fixed spatial origin]
\label{lem:stationary-zero-hold}
Under the stationary environment law \(\cF_\lambda\), conditional on
\(K(0)=0\), let \(A_0\) be the time since the last return to zero and
\(B_0\) the time until the next departure.  Then \(A_0,B_0\) are independent
\(\operatorname{Exp}(2)\) variables.  Reindex the cells so this zero hold
belongs to cell zero.  The entire marked excursion field and all other
zero holds have their ordinary return-Palm joint law, independently of
\((A_0,B_0)\).
\end{lemma}

\begin{proof}
Let \(\rho_{\mathcal R}>0\) be the return intensity.  The intervals
\([R_j,R_{j+1})\) partition the line.  Campbell's formula applied to this
partition gives the stationary renewal inversion
\begin{equation}
\E F(\omega)
=\rho_{\mathcal R}\E^{\mathcal R}
  \int_0^{L_0}F(\theta_t\omega)\,dt
\label{eq:8.7}
\end{equation}
for every nonnegative measurable environment function \(F\).  Indeed,
integrate \(F(\theta_x\omega)\) over \([0,1]\), split the integral over
the return intervals, and use stationarity and Campbell; Tonelli removes
the spatial unit-window indicator.  In particular,
\(\rho_{\mathcal R}\E^{\mathcal R}L_0=1\).

Write \(\Xi\) for the marked excursions and the holds \((Z_j)_{j\ne0}\),
all indexed relative to the return starting the distinguished cell.
Lemma~\ref{lem:return-cells-prove} makes \(\Xi\) independent of \(Z_0\).
Except at event coordinates, \(K(t)=0\) in cell zero exactly when
\(0<t<Z_0\).  For any nonnegative measurable \(\phi\), inversion and
the density \(2e^{-2z}\) of \(Z_0\) therefore give
\begin{equation}
\begin{aligned}
\E\bigl[\mathbf1_{\{K(0)=0\}}\phi(A_0,B_0,\Xi)\bigr]
&=\rho_{\mathcal R}\E^{\mathcal R}
    \int_0^{Z_0}\phi(t,Z_0-t,\Xi)\,dt\\
&=\rho_{\mathcal R}\int_0^\infty\!\int_0^\infty
    2e^{-2(a+b)}\E^{\mathcal R}\phi(a,b,\Xi)\,da\,db.
\end{aligned}
\label{eq:8.8}
\end{equation}
Here \(\Xi\) on the left is read after reindexing at the last return;
it does not depend on the position within the zero hold.  Taking
\(\phi=1\) yields
\(\Pp\{K(0)=0\}=\rho_{\mathcal R}/2=\varpi(0)>0\).
Dividing \eqref{eq:8.8} by this probability gives the product density
\(4e^{-2a}e^{-2b}\) times the original law of \(\Xi\), proving the
full joint assertion for the environment law.
\end{proof}

\subsection{Conditional environmental regularity}

For \(h\ge1\), define an event in the original spatial stationary law by
\begin{equation}
\mathsf A_h:=\{K(0)=0,\ h\le B_0\le h+1\},\qquad
p_h:=\Pp(\mathsf A_h)
=\varpi(0)e^{-2h}(1-e^{-2})>0.
\label{eq:8.9}
\end{equation}
For any interval \(I\) of length \(L\) and center \(c\), let \(I^*\)
be its concentric length-\(4L\) enlargement and set
\[
J^-=[c-2L,c-L],\qquad J^+=[c+L,c+2L].
\]
Given constants \(C_0,c_0>0\), let \(\mathsf R(I)\) mean that
\begin{enumerate}[label=\textup{(\roman*)},leftmargin=2.3em]
\item between every pair of non-event zero-current cuts
\(x^-\in J^-\), \(x^+\in J^+\), there are at most \(C_0L\)
source events;
\item the sum of \(V_j\) over the complete return-cell event sets between
any such pair is at most \(C_0L\);
\item each of \(J^-,J^+\) has zero-current occupation at least \(c_0L\).
\end{enumerate}
The counts of the two colours between these cuts agree, because those
events form a union of complete return-cell event sets.

\begin{lemma}[Environmental regularity conditional on a long residual hold]
\label{lem:conditional-gap-regularity}
There are constants \(C_0,c_0,C,L_*\), depending only on the feedback law,
such that, if \(I\) has center \(h/2\) and length
\(L\ge\max\{L_*,16(h+1)\}\), then
\begin{equation}
\Pp\{\mathsf R(I)^c\mid\mathsf A_h\}\le \frac{C}{L}.
\label{eq:8.10}
\end{equation}
The constants do not depend on \(h\) or on any matching phase.
\end{lemma}

\begin{proof}
By Lemma~\ref{lem:stationary-zero-hold}, conditional on \(\mathsf A_h\),
the age \(A_0\) remains a rate-two exponential, independent of \(B_0\)
and of the ordinary excursion and neighboring-cell data.

Collapse \([0,B_0]\), leaving the negative half-line fixed and translating the
subsequent environment left by \(B_0\). In the auxiliary environment the
distinguished return is at \(-A_0\), its departure is at zero, and its
next return is at \(E_0\). The distinguished cell now has zero hold
\(A_0\), so the entire indexed cell array has the ordinary iid Palm law,
independently of \(h\) and \(B_0\). The original matching and its
conditional law are retained throughout.

It suffices to control the partial sums of four nonnegative cell
observables: length \(L_j\), zero hold \(Z_j\), event count \(N_j\),
and reward \(V_j\). Write \(\mu_X=\E X_0\) for any of these
observables. All have finite second moments, and \(\mu_L,\mu_Z>0\).
Put \(k=\lfloor L/M\rfloor\). For each of the two index directions,
let \(S_X(n)\) be the sum of the first \(n\) cell observables in that
direction. For fixed integers \(M,H\), Chebyshev's inequality and a
finite union bound give
\[
\Pp\bigl\{\exists X,m:\ |S_X(mk)-\mu_Xmk|>k,
\ 0\le m\le H\bigr\}\le C_H/k.
\]
On the complementary event, monotonicity between consecutive grid
indices gives
\[
|S_X(n)-\mu_Xn|\le(1+\mu_X)k,
\qquad 0\le n\le Hk,
\]
in both directions. This interpolation constant is independent of
\(H\).

Differences of these bounds also control individual boundary
cells by \(Ck\). In particular \(A_0\le L_0\le Ck\), so the
auxiliary return coordinates differ from their linear means by at most
\(Ck\).

Choose \(M\) first, large enough that these errors are a small
fraction of \(L\), and then choose \(H\), depending on \(M\), so
that the controlled returns extend beyond \([-3L,3L]\). Increase
\(L_*\) so that \(k\ge L/(2M)\). The event-count and reward sums
then bound the total counts and rewards of cells meeting this window
by \(CL\).

For a subinterval of length \(\ell\ge L/2\) in the
window, discarding its two boundary cells loses at most \(Ck\) in
length. Subtracting the length and zero-hold partial sums over the
remaining complete cells gives zero occupation at least
\[
\frac{\mu_Z}{\mu_L}\ell-Ck\ge c_0L,
\]
after the preceding choice of \(M\), with constants depending only
on the cell law. Event coordinates have zero Lebesgue measure. The
exceptional probability is at most \(C_H/k\le C'/L\).

To transfer these estimates back, put \(c=h/2\). Then
\(c\le L/32\) and \(B_0\le L/16\). A pair of original search
cuts corresponds in the auxiliary environment to
\(x^-\) and \(x^+-B_0\), both lying in \([-3L,3L]\).
The event sets and their cell rewards are unchanged, so the preceding
upper bounds give (i) and (ii).

For the lower occupation bounds, use the fixed auxiliary intervals
\[
K^-=J^-,\qquad
K^+=[c+L,c+2L-L/16].
\]
Both lie in \([-3L,3L]\) and have length at least \(L/2\).
The negative interval is unchanged when the hold is restored, while
\(K^++B_0\subset J^+\) for every \(0\le B_0\le L/16\).
Translation preserves its zero occupation, proving (iii) in the
original environment. Thus all three estimates are uniform over the
conditioned values \(B_0\in[h,h+1]\).
\end{proof}

\subsection{The deterministic energy iteration}

\begin{lemma}[Zero-current surgery]
\label{lem:zero-surgery-prove}
Let \(x\) be a cut which is not an event coordinate, with \(K(x)=0\), and suppose
\begin{equation}
T_\Sigma(x)\le\Theta,
\qquad
S_\Sigma(x)\le\Theta.
\label{eq:8.11}
\end{equation}
Delete the selected edges crossing \(x\).  The unmatched vertices on either side of the
cut can be paired separately on that side, and the total feedback cost of all new repair
edges is at most
\begin{equation}
C_\lambda\Theta^3.
\label{eq:8.12}
\end{equation}
\end{lemma}

\begin{proof}
Zero relative flux gives oriented crossing current \(K(x)=0\), so the unmatched
vertices on each side have equal source and target counts.  There are at most
\(T_\Sigma(x)\le\Theta\) crossing edges.

Every endpoint of a crossing edge is at distance at most the length of
that edge from the cut, and every individual crossing
length is at most \(S_\Sigma(x)\le\Theta\).  Hence all unmatched endpoints on a fixed
side lie within distance \(\Theta\) of the cut.  Pairing them arbitrarily produces at
most \(\Theta\) new edges, each of length at most \(2\Theta\).  The upper coercivity
bound in \eqref{eq:3.11} gives \eqref{eq:8.12}.
\end{proof}

\begin{lemma}[One-step bootstrap]
\label{lem:one-step-bootstrap-prove}
For fixed \(C_0,c_0>0\), there are constants \(C_B,r_*,L_*\) such that
the following holds. Let \(L\ge L_*\), \(r\ge r_*\), and choose
\(\Theta>0\) satisfying
\begin{equation}
\Theta<L/4,\qquad C_B\Theta^3\le rL/2.
\label{eq:8.13}
\end{equation}
On the regular environment event \(\mathsf R(I)\),
\begin{equation}
\mathcal H_\Sigma(I)>rL
\quad\Longrightarrow\quad
\mathcal H_\Sigma(I^*)>\frac{c_0}{8}\Theta\,|I^*|.
\label{eq:8.14}
\end{equation}
\end{lemma}

\begin{proof}
Suppose first that each search interval contains a zero-current cut \(x^-\in J^-\),
\(x^+\in J^+\), chosen away from the countable event coordinates, satisfying
\begin{equation}
T_\Sigma(x^\pm)\le\Theta,
\qquad
S_\Sigma(x^\pm)\le\Theta.
\label{eq:8.15}
\end{equation}
Because the cuts are at distance at least \(L/2>2\Theta\) from \(I\), an edge crossing
one of the cuts cannot meet \(I\).  Also no edge can cross both cuts.

Delete all selected edges meeting either cut and every selected edge with
both endpoints between the cuts.  Match all source and target events
between the cuts by the cellwise increasing environmental matching.
Each non-event zero cut lies in an
initial zero hold, so the events between the cuts are a union of complete
return-cell event sets \((R_j,R_{j+1}]\), each with equal colour counts.
Repair the unmatched exterior endpoints separately at the two cuts using
Lemma~\ref{lem:zero-surgery-prove}.  The resulting perfect matching differs from
\(\Sigma\) on finitely many edges.

Local minimality, the lower coercivity in \eqref{eq:3.11}, the environmental reward bound, and
the two repair estimates give
\begin{equation}
\sum_{e:I_e\cap I\ne\varnothing}d_e^2
\le C\bigl(L+\Theta^3\bigr).
\label{eq:8.16}
\end{equation}
At most \(CL\) selected edges meet \(I\), since their source endpoints
lie between the cuts.  Using \(|d|\le(1+d^2)/2\), we obtain
\begin{equation}
\mathcal H_\Sigma(I)
\le\sum_{e:I_e\cap I\ne\varnothing}(|d_e|+d_e^2)
\le \frac{CL}{2}
+\frac32\sum_{e:I_e\cap I\ne\varnothing}d_e^2
\le C'(L+\Theta^3).
\label{eq:8.17}
\end{equation}
Take \(C_B\ge C'\) and \(r_*\ge2C_B\). The cubic condition gives
\begin{equation}
\mathcal H_\Sigma(I)\le C_BL+C_B\Theta^3\le rL.
\label{eq:8.18}
\end{equation}
Thus \(\mathcal H_\Sigma(I)>rL\) is impossible if both good cuts exist.

Hence on that event at least one search interval contains no zero-current cut satisfying
\eqref{eq:8.15}.  On the regular environment event, the zero set in that search interval has
Lebesgue measure at least \(c_0L\).  At almost every point of that zero set,
\(T_\Sigma+S_\Sigma>\Theta\).  Therefore
\[
\mathcal H_\Sigma(I^*)\ge c_0\Theta L
=\frac{c_0}{4}\Theta|I^*|.
\]
This is strictly larger than \((c_0/8)\Theta|I^*|\), proving \eqref{eq:8.14}.
\end{proof}

\begin{theorem}[Local energy conditional on a long zero hold]
\label{thm:conditional-gap-energy}
There are constants \(A,C,L_*\) such that, for every
\(\Sigma\in\mathcal G_{\lambda,2}^0\), every \(h\ge1\), and
every \(L\ge\max\{L_*,16(h+1)\}\), if \(I\) has center
\(h/2\) and length \(L\), then
\begin{equation}
\Pp\left\{
\mathcal H_\Sigma(I)>AL\mid\mathsf A_h\right\}\le \frac{C}{L}.
\label{eq:8.19}
\end{equation}
\end{theorem}

\begin{proof}
Use the constants from Lemmas~\ref{lem:conditional-gap-regularity}
and~\ref{lem:one-step-bootstrap-prove}. Choose \(A\ge r_*\) and
increase \(L_*\) so that, for every \(L\ge L_*\),
\[
\frac{16A}{c_0}<\frac L4,
\qquad
C_B\left(\frac{16A}{c_0}\right)^3\le\frac{AL}{2}.
\]
Let \(I_n\) be concentric intervals centered at \(h/2\), with
\[
L_n=4^nL,\qquad r_n=2^nA,\qquad
\Theta_n=\frac{16r_n}{c_0}.
\]
The cubic inequality persists at every scale, since \(r_n^2/L_n=A^2/L\);
also \(\Theta_n/L_n\) decreases. The one-step conclusion has exactly
the next threshold, because \((c_0/8)\Theta_n=r_{n+1}\). Hence
\[
\{\mathcal H_\Sigma(I_n)>r_nL_n\}
\subset
\{\mathcal H_\Sigma(I_{n+1})>r_{n+1}L_{n+1}\}
\cup\{\mathsf R(I_n)^c\}.
\]
This deterministic inclusion holds under the conditional joint law.
Iteration gives
\begin{equation}
\begin{aligned}
\Pp\{\mathcal H_\Sigma(I_0)>AL\mid\mathsf A_h\}
\le{}&\Pp\{\mathcal H_\Sigma(I_N)>r_NL_N\mid\mathsf A_h\}\\
&+C\sum_{n=0}^{N-1}\frac1{4^nL}.
\end{aligned}
\label{eq:8.20}
\end{equation}
By Lemma~\ref{lem:mean-local-energy-prove} and Markov's inequality,
the first term is at most
\[
\frac{C_\Sigma}{p_h r_N}\longrightarrow0.
\]
Here \(h\) is fixed, so \(p_h>0\) is fixed. The remaining series is
\(4/(3L)\), proving \eqref{eq:8.19} after changing \(C\). Thus the conditioning
factor \(p_h^{-1}\) affects only the vanishing terminal term.
\end{proof}

\begin{corollary}[A common empty cut in any stationary joining]
\label{cor:common-empty-gap}
Let \((\Sigma^+,\Sigma^-)\) be a stationary common-environment
joining of two elements of \(\mathcal G_{\lambda,2}^0\).  In its
original spatial stationary law, there is positive probability of a
cut crossed by neither matching.
\end{corollary}

\begin{proof}
Take \(L=16(h+1)\) and \(I\) centered at \(h/2\).  On
\(\mathsf A_h\), there is no event in \((0,h)\).  Any edge
crossing \(h/2\) therefore has length at least \(h\), crosses
all of \((0,h)\), and contributes at least \(h^2\) to
\(\int_0^h S_\Sigma(x)\,dx\).

For sufficiently large fixed
\(h\), one has \(h^2>AL\) and the sum of the two bounds in
\eqref{eq:8.19} is less than \(1/2\).  Conditional on \(\mathsf A_h\),
the union bound then gives probability at least \(1/2\) that neither
matching crosses \(h/2\).  Its unconditional spatial probability
is at least \(p_h/2>0\).
\end{proof}

\subsection{Disagreement components and uniqueness}

\begin{lemma}[Every stationary disagreement component spans both infinities]
\label{lem:spanning-disagreements}
In a stationary common-environment joining of two elements of
\(\mathcal G_{\lambda,2}^0\), almost surely every component of the
symmetric difference has one end tending to \(-\infty\) and one
end tending to \(+\infty\).  Consequently, a common empty cut
implies equality of the entire matchings.
\end{lemma}

\begin{proof}
After common edges are removed, every remaining vertex has degree two.
Thus each component is a finite alternating cycle or a double ray.
The first possibility is excluded by local minimality and
Lemma~\ref{lem:no-finite-alternating-ties}: switching the cycle in
each matching would otherwise force its two costs to be equal.

Lemma~\ref{lem:finite-crossings} holds simultaneously at all rational
cuts.  Every end of a double ray therefore tends to one of the two
spatial infinities.  Otherwise local finiteness would fail in a bounded
interval, or that ray would traverse a rational cut infinitely often.

If both ends tended right, the component would have a unique leftmost
vertex.  Send one unit of mass from each vertex of every such component
to its leftmost vertex.  This transport is translation covariant and
has outgoing intensity at most the finite total vertex intensity.
Each recipient, however, receives infinite mass.  The stationary
mass-transport identity excludes every such recipient almost surely.
The analogous transport to a rightmost vertex excludes two left ends.

Every remaining component thus crosses every cut.  If both matchings
have no edge crossing one cut, no disagreement component can exist.
\end{proof}

\begin{theorem}[Regenerative uniqueness at finite energy]
\label{thm:regenerative-matching-uniqueness}
Any stationary common-environment joining of two matchings in
\(\mathcal G_{\lambda,2}^0\) is supported on equal matchings.
\end{theorem}

\begin{proof}
Suppose the invariant event
\(D=\{\Sigma^+\ne\Sigma^-\}\) has positive probability, and
condition the stationary joining on \(D\).  The feedback environment
is ergodic, as follows from its iid return cells.  Conditioning an
extension over an ergodic environment on a translation-invariant event
leaves its environmental marginal unchanged.  Indeed, for a bounded
environmental test function \(f\), invariance gives
\[
\E[\mathbf1_D f]
=\E\left[\mathbf1_D\,
\frac1{2T}\int_{-T}^T f\circ\theta_x\,dx\right]
\longrightarrow \Pp(D)\E f
\]
by the environmental ergodic theorem and bounded convergence.

The conditioned matching marginals remain stationary, locally
minimizing and of zero relative flux.  Their source-Palm second moments
remain finite: the corresponding unnormalized stationary Palm
integrals are bounded by the original integrals divided by \(\Pp(D)\),
and the environmental source intensity is unchanged.

Corollary~\ref{cor:common-empty-gap} therefore applies to this
conditioned joining and gives positive spatial probability of a
common empty cut.  Lemma~\ref{lem:spanning-disagreements} then forces
the matchings to coincide, contradicting \(D\).  Thus \(\Pp(D)=0\).
\end{proof}

\begin{corollary}[A unique matching kernel over the environment]
\label{cor:unique-environment-kernel}
There is at most one stationary matching law in
\(\mathcal G_{\lambda,2}^0\).  When such a law exists, its
conditional matching kernel given the feedback environment is almost
surely a point mass.
\end{corollary}

\begin{proof}
Join any two regular conditional matching kernels independently over
their common feedback environment.  Stationarity of the joint laws
implies stationarity of this relative independent joining.  Theorem
\ref{thm:regenerative-matching-uniqueness} makes the two samples equal
almost surely.  Applying the same assertion to two conditional copies
of one law shows that its kernel is a point mass.  Applying it to two
different laws shows that their kernels coincide almost surely.
\end{proof}

\begin{theorem}[Projective uniqueness]
\label{thm:EFE2}
Every \(\Sigma\in\mathcal G_{\lambda,2}^0\) has the same source-Palm
projective law \(\kappa[\Sigma]=\kappa_\infty\).
\end{theorem}

\begin{proof}
The class is nonempty by the good-event compactness theorem.
Corollary~\ref{cor:unique-environment-kernel} identifies its stationary
matching law, and hence its source-Palm projective law.  Its first
mark marginal is uniform by the source intensity \eqref{eq:3.9}.

For the second, send \(\mathbf1_B(v)\) along every matched edge to its target.
Stationary mass transport and perfectness identify the source-window
outgoing mass with the target-window count carrying marks in \(B\).
Its expectation is \(|B|\) by the target intensity \eqref{eq:3.9}.  Thus the
second marginal is also uniform.
\end{proof}

\section{Macroscopic reconstruction}
\label{sec:reconstruction}

\subsection{Vertical concentration}

\begin{proposition}[Vertical diagonal concentration]
\label{prop:vd}
For the Euclidean optimizer,
\begin{equation}
\boxed{
\frac1n\sum_{i=1}^n|S_{\sigma_n(i)}-T_i|
\xrightarrow{\Pp}0.
}
\label{eq:5.1}
\end{equation}
\end{proposition}

\begin{proof}
For a permutation \(\sigma\), set
\[
\mathcal E_n(\sigma)
=
\frac1n\sum_i
D_{i,\sigma(i)}
L\left(\frac{S_{\sigma(i)}-T_i}{D_{i,\sigma(i)}}\right).
\]
The Euclidean objective differs from \(\mathcal E_n\) by the
permutation-independent average of \(D_{ij}=2+V_j-U_i\).  Let \(\rho_n\)
match the order statistics of \(T\) and \(S\).  By DKW,
\[
\max_i|T_{(i)}-S_{(i)}|=O_{\Pp}(n^{-1/2}).
\]
Since \(L(z)\le z^2/2\),
\[
\mathcal E_n(\sigma_n)
\le \mathcal E_n(\rho_n)
\le \frac1{2n}\sum_i|T_{(i)}-S_{(i)}|^2
=O_{\Pp}(n^{-1}).
\]
For \(|z|\le1\),
\[
L(z)\ge \frac{z^2}{1+\sqrt2}.
\]
Because \(1\le D_{ij}\le3\),
\[
\mathcal E_n(\sigma_n)
\ge
\frac{1}{3(1+\sqrt2)}
\frac1n\sum_i|S_{\sigma_n(i)}-T_i|^2.
\]
Hence the average squared vertical displacement is \(O_{\Pp}(n^{-1})\),
and Cauchy--Schwarz gives \eqref{eq:5.1}.
\end{proof}

\subsection{Block identification and reconstruction}

Let \(\kappa_\infty\) be the common projective law from Theorem~\ref{thm:EFE2}.

\begin{lemma}[Macroscopic block identification]
\label{lem:block-comp}
Let \(J=[a,b]\Subset(0,1)\) be fixed.
For every bounded continuous \(f:[0,1]^2\to\R\),
\begin{equation}
\frac1n\sum_{i=1}^n
\mathbf 1_{\{T_i\in J\}}
f(U_i,V_{\sigma_n(i)})
\xrightarrow{\Pp}
|J|\int f\,d\kappa_\infty .
\label{eq:6.1}
\end{equation}
\end{lemma}

\begin{proof}
Fix \(\varepsilon>0\) and use its original feedback graph.  Suppose
for some \(\delta>0\) the events
\[
\mathcal G_{n,\varepsilon}\cap
\left\{\left|\frac1n\sum_{\widehat\xi_i\in nJ}
 f(U_i,V_{\sigma_n(i)})-|J|\int f\,d\kappa_\infty\right|>\delta\right\}
\]
have probability bounded below along a subsequence.  Apply
Lemma~\ref{lem:good-sample-extraction} within these events, now with
roots in \(nJ\).  The same stopped estimates apply: all fixed windows
lie inside the bulk eventually, and the global current and energy
bounds acquire the fixed factor \(|J|^{-1}\).

The proof of Theorem~\ref{thm:comp} gives a further deterministic weak limit
\(Q\in\mathcal G_{\lambda,2}^0\).  Corollary~\ref{cor:I4} and the
global VC event give the two count inputs of Lemma~\ref{lem:palm}.
Its Palm readout and Theorem~\ref{thm:EFE2} show that the displayed block statistic tends
to \(|J|\int f\,d\kappa_\infty\), a contradiction.
Thus its good-event failure probability tends to zero.

To replace intersection coordinates by physical heights, put
\(\zeta_{n,\varepsilon}=\|g_{n,\varepsilon}\|_\infty/n\) and
\(\Delta_n=\sup_t|F_n(t)-t|\).
The symmetric difference between
\(\{i:\widehat\xi_i/n\in J\}\) and \(\{i:T_i\in J\}\)
is contained in height strips of width \(\zeta_{n,\varepsilon}\)
around the two endpoints.  Its normalized size is at most
\begin{equation}
4\zeta_{n,\varepsilon}+4\Delta_n+O(1/n)
 =o_{\Pp}(1)\quad\hbox{on }\mathcal G_{n,\varepsilon}.
\label{eq:6.4d}
\end{equation}
Multiplication by \(\|f\|_\infty\) bounds the statistic's change.
The physical block failure probability therefore has limsup at most
\(\varepsilon\).  This statistic and its asserted limit are independent
of the auxiliary tolerance; letting \(\varepsilon\downarrow0\)
proves convergence in probability.
\end{proof}

\begin{theorem}[Blockwise reconstruction of the empirical plan]
\label{thm:reconstruct}
The empirical plans satisfy
\begin{equation}
\boxed{
\pi_n
\xrightarrow{\Pp}
dt\,\delta_t(ds)\,\kappa_\infty(du,dv).
}
\label{eq:6.4}
\end{equation}
\end{theorem}

\begin{proof}
By Lemma~\ref{lem:block-comp}, for every interval
\(J\Subset(0,1)\) with rational endpoints and every bounded continuous
\(f(u,v)\),
\begin{equation}
\frac1n\sum_i
\mathbf 1_{\{T_i\in J\}}f(U_i,V_{\sigma_n(i)})
\xrightarrow{\Pp}
|J|\int f\,d\kappa_\infty .
\label{eq:6.5}
\end{equation}
Finite linear combinations of such interval indicators approximate every
continuous \(\phi:[0,1]\to\R\) uniformly after controlling arbitrarily
small endpoint intervals.  Hence
\begin{equation}
\frac1n\sum_i
\phi(T_i)f(U_i,V_{\sigma_n(i)})
\xrightarrow{\Pp}
\left(\int_0^1\phi(t)\,dt\right)
\left(\int f\,d\kappa_\infty\right).
\label{eq:6.6}
\end{equation}
By Stone--Weierstrass, finite sums of products
\(\phi(t)f(u,v)\) are uniformly dense in
\(C([0,1]^3)\).  Therefore
\begin{equation}
\frac1n\sum_i
\delta_{(T_i,U_i,V_{\sigma_n(i)})}
\xrightarrow{\Pp}
dt\,\kappa_\infty(du,dv).
\label{eq:6.7}
\end{equation}

Finally, for any Lipschitz test
\(F(t,u,s,v)\),
\[
\left|
\frac1n\sum_i
F(T_i,U_i,S_{\sigma_n(i)},V_{\sigma_n(i)})
-
\frac1n\sum_i
F(T_i,U_i,T_i,V_{\sigma_n(i)})
\right|
\le
\Lip(F)\,
\frac1n\sum_i|S_{\sigma_n(i)}-T_i|,
\]
which tends to zero in probability by Proposition~\ref{prop:vd}.  Approximation of continuous
tests by Lipschitz tests on the compact state space then gives \eqref{eq:6.4}.
\end{proof}

\begin{remark}[The feedback parameter is auxiliary]
The empirical plans in Theorem~\ref{thm:reconstruct} depend only on the
original samples.  Running the argument with two admissible feedback
parameters therefore gives two weak limits of the same sequence, which
must coincide.  Thus \(\kappa_\infty\) is independent of \(\lambda\).
\end{remark}

\subsection{Conclusion of the selection proof}

\begin{proof}[Proof of Theorem~\ref{thm:main-prove}]
Theorem~\ref{thm:EFE2} supplies the unique source-Palm horizontal coupling.
Theorem~\ref{thm:reconstruct} combines physical block identification and
vertical concentration to give the asserted convergence.
The two marginals of \(\kappa_\infty\)
are uniform because the finite empirical mark marginals are the laws
of all \(U_i\)'s and all \(V_j\)'s, respectively.
\end{proof}
\begin{corollary}[Deterministic projective mark]
\label{cor:deterministic-mark-prove}
The original full empirical mark coupling
\(\kappa_n=n^{-1}\sum_i\delta_{(U_i,V_{\sigma_n(i)})}\)
converges weakly in probability to \(\kappa_\infty\).
\end{corollary}

\begin{proof}
Apply Theorem~\ref{thm:main-prove} to tests of the form
\(F(t,u,s,v)=f(u,v)\), where \(f\) is continuous on \([0,1]^2\).
\end{proof}

\begin{thebibliography}{9}

\bibitem{AmbrosioPratelli}
L.~Ambrosio and A.~Pratelli.
Existence and stability results in the \(L^1\) theory of optimal transportation.
In \emph{Optimal Transportation and Applications},
Lecture Notes in Mathematics, vol.~1813, pp.~123--160.
Springer, 2003.
\href{https://doi.org/10.1007/978-3-540-44857-0_5}{doi:10.1007/978-3-540-44857-0\_5}.

\bibitem{delBarrioEtAl2025}
E.~del Barrio, A.~Gonz\'alez-Sanz, J.-M.~Loubes, and
D.~Rodr\'iguez-V\'itores.
Distributional limit theory for optimal transport.
Preprint, 2025, revised 2026.
\href{https://arxiv.org/abs/2505.19104}{arXiv:2505.19104}.

\bibitem{Gwozdz}
M.~Gw\'{o}\'zd\'z.
Instability of the ray-monotone selector for \(W_1\)-optimal transport.
Preprint, 2026.
\href{https://arxiv.org/abs/2604.15474}{arXiv:2604.15474}.

\bibitem{HolroydJansonWastlund}
A.~E.~Holroyd, S.~Janson, and J.~W\"astlund.
Minimal matchings of point processes.
\emph{Probability Theory and Related Fields} \textbf{184} (2022), 571--611.
\href{https://doi.org/10.1007/s00440-022-01151-y}{doi:10.1007/s00440-022-01151-y}.

\bibitem{Huesmann}
M.~Huesmann.
Optimal transport between random measures.
\emph{Annales de l'Institut Henri Poincar\'e, Probabilit\'es et Statistiques}
\textbf{52} (2016), 196--232.
\href{https://doi.org/10.1214/14-AIHP634}{doi:10.1214/14-AIHP634}.

\end{thebibliography}
\end{document}
